%% 4-5 面積と体積：細かく分けてたす（第13回：新しく書いた。組 V4-3。ページ制限なし。偶数ページ）
%% ほかの巻が指す形：
%%   5-9「三角形の面積＝底辺×高さ÷2」「台形＝（上底＋下底）×高さ÷2（同じ台形を2つ、さかさに合わせると平行四辺形）」
%%   5-10「円の面積は、細いおうぎ形に切って並べかえ、長方形に近づけて求めた」「円板（柱）の体積＝底面積×厚さ」「円すい 1/3πr²h」
%%   5-4・5-6「球の体積 4/3πr³、表面積 4πr²、円の面積 πr²、円すいの体積＝底面積×高さ×1/3」
%%   三角形の言い方（第13回）：直角三角形は2つ合わせると長方形。ふつうの三角形は2つ合わせると平行四辺形（それを切って移すと長方形）
\makeatletter
% 本文の中の小見出し（それだけで1行）  \bkivfhead{見出し}
\newcommand\bkivfhead[1]{\par\addvspace{4pt}\Needspace{4\baselineskip}\noindent{\sffamily\bfseries #1}\par\nobreak\noindent\ignorespaces}
% 式を、まん中に1行で置く  \bkivfline{…}
\newcommand\bkivfline[1]{\par\addvspace{3pt}{\centering #1\par}\addvspace{3pt}\noindent\ignorespaces}
% 説明の段の見出し（丸数字＋太字の1行）  \bkivfstep{番号}{見出し}
\newcommand\bkivfstep[2]{\par\addvspace{6pt}\Needspace{5\baselineskip}\noindent
  {\bkcircled[bkOrange]{#1}\enspace\sffamily\bfseries\color{bkOrange!75!black}#2}\par\nobreak\vspace{2pt}\noindent\ignorespaces}
% 枠でかこんだ大事な式  \bkivfwaku{…}
\newcommand\bkivfwaku[1]{\par\addvspace{4pt}{\centering\setlength\fboxsep{5pt}\fcolorbox{bkOrange}{white}{#1}\par}\addvspace{4pt}\noindent\ignorespaces}
% 直角のしるし  \bkivfright{(頂点)}{(一方の点)}{(もう一方の点)}
\newcommand\bkivfright[3]{%
  \coordinate (bkivfp) at ($#1!2mm!#2$);\coordinate (bkivfq) at ($#1!2mm!#3$);
  \draw[bkInk, line width=0.5pt] (bkivfp) -- ($(bkivfp)+(bkivfq)-#1$) -- (bkivfq);}
% 円を n 等分したおうぎ形（色を交互に）  \bkivfpie{(中心)}{n}{半径 mm}
\newcommand\bkivfpie[3]{\begin{scope}[shift={#1}]
  \pgfmathsetmacro\bkivfA{360/#2}%
  \foreach \k in {1,...,#2} {\pgfmathsetmacro\bkivfa{(\k-1)*\bkivfA}\pgfmathsetmacro\bkivfb{\k*\bkivfA}%
    \pgfmathparse{int(mod(\k,2))}\ifnum\pgfmathresult=1\def\bkivfc{bkOrange!35}\else\def\bkivfc{bkBlue!25}\fi
    \filldraw[fill=\bkivfc, draw=bkInk!70, line width=0.4pt] (0,0) -- (\bkivfa:#3) arc (\bkivfa:\bkivfb:#3) -- cycle;}
\end{scope}}
% n 等分したおうぎ形を、上向き・下向き交互に並べる  \bkivfrearr{(左下)}{n}{半径 mm}
\newcommand\bkivfrearr[3]{\begin{scope}[shift={#1}]
  \pgfmathsetmacro\bkivfh{180/#2}\pgfmathsetmacro\bkivfw{2*#3*sin(\bkivfh)}%
  \pgfmathsetmacro\bkivfm{int(#2/2-1)}%
  \foreach \k in {0,...,\bkivfm} {%
    \filldraw[fill=bkOrange!35, draw=bkInk!70, line width=0.4pt] (\k*\bkivfw,0) -- ++(90-\bkivfh:#3) arc (90-\bkivfh:90+\bkivfh:#3) -- cycle;
    \filldraw[fill=bkBlue!25, draw=bkInk!70, line width=0.4pt] ({#3*sin(\bkivfh)+\k*\bkivfw},{#3*cos(\bkivfh)}) -- ++(270-\bkivfh:#3) arc (270-\bkivfh:270+\bkivfh:#3) -- cycle;}
\end{scope}}
% 斜めから見た箱  \bkivfbox{x0}{y0}{z0}{dx}{dy}{dz}{色}（3-8 の \bkthribox と同じぬり方）
\newcommand\bkivfbox[7]{%
  \pgfmathsetmacro\bkivfX{#1+#4}\pgfmathsetmacro\bkivfY{#2+#5}\pgfmathsetmacro\bkivfZ{#3+#6}%
  \filldraw[fill=#7!42, draw=#7!60!black, line width=0.5pt, line join=round] (#1,#2,#3) -- (\bkivfX,#2,#3) -- (\bkivfX,#2,\bkivfZ) -- (#1,#2,\bkivfZ) -- cycle;
  \filldraw[fill=#7!18, draw=#7!60!black, line width=0.5pt, line join=round] (#1,#2,\bkivfZ) -- (\bkivfX,#2,\bkivfZ) -- (\bkivfX,\bkivfY,\bkivfZ) -- (#1,\bkivfY,\bkivfZ) -- cycle;
  \filldraw[fill=#7!68, draw=#7!60!black, line width=0.5pt, line join=round] (\bkivfX,#2,#3) -- (\bkivfX,\bkivfY,#3) -- (\bkivfX,\bkivfY,\bkivfZ) -- (\bkivfX,#2,\bkivfZ) -- cycle;}
% 横から見た円柱（中心の下 (x,y)、半径 r、高さ h、色）  \bkivfcyl{(x,y)}{r}{h}{色}
\newcommand\bkivfcyl[4]{\begin{scope}[shift={#1}]
  \fill[#4!22] (-#2,0) -- (-#2,#3) arc (180:360:#2 and #2*0.3) -- (#2,0) arc (360:180:#2 and #2*0.3) -- cycle;
  \filldraw[fill=#4!10, draw=#4!60!black, line width=0.6pt] (0,#3) ellipse (#2 and #2*0.3);
  \draw[#4!60!black, line width=0.6pt] (-#2,#3) -- (-#2,0) arc (180:360:#2 and #2*0.3) -- (#2,#3);
  \draw[#4!60!black, line width=0.4pt, dashed] (-#2,0) arc (180:0:#2 and #2*0.3);
\end{scope}}
% 横から見た円すい（底の中心 (x,y)、半径 r、高さ h、色）  \bkivfcone{(x,y)}{r}{h}{色}
\newcommand\bkivfcone[4]{\begin{scope}[shift={#1}]
  \fill[#4!22] (-#2,0) -- (0,#3) -- (#2,0) arc (360:180:#2 and #2*0.3) -- cycle;
  \draw[#4!60!black, line width=0.6pt] (-#2,0) -- (0,#3) -- (#2,0) arc (360:180:#2 and #2*0.3);
  \draw[#4!60!black, line width=0.4pt, dashed] (-#2,0) arc (180:0:#2 and #2*0.3);
\end{scope}}
\makeatother

\bkchapter{4}{5}{面積と体積：\\細かく分けてたす}{\bklevelrange{1}{2}}{面積と体積：細かく分けてたす}
\bkmichishirube{1-21,2-7,2-9,2-10,4-2,4-3,4-4}{5-4,5-6,5-9,5-10}

\begin{goal}
\begin{itemize}[label={\color{\bkTT}●}, leftmargin=1.6\zw, itemsep=0pt, topsep=1pt]
\item \bklevel{1}\enspace 平行四辺形・三角形・台形・ひし形の面積を、長方形に直す考え方で求められる。
\item \bklevel{2}\enspace 円の面積 $\pi r^2$ と、柱・すい・球の体積を、細かく分けて考えたわけといっしょに使える。
\item \bklevel{2}\enspace 展開図から表面積を求め、体積を単位をつけて L などに換算できる。
\end{itemize}
\end{goal}

\begin{hatena}
長方形の面積は「たて×横」で求められました（1-21）。では、円の面積はどうでしょう。円には、まっすぐな辺がありません。

そこで、円をピザのように8つに切って、とがった所を上と下に交互に向けて並べてみます。すると、でこぼこしているけれど、長方形に少しにた形になりました。もっと細かく切ったら、どうなるでしょう。
\begin{bkfigure}
\begin{tikzpicture}[x=1mm, y=1mm]
  \bkivfpie{(16,14)}{8}{16}
  \draw[bk arrow, bkGray, line width=0.9pt] (36,14) -- (46,14);
  \bkivfrearr{(50,6.6)}{8}{16}
  \node[font=\sffamily\bfseries\Large, text=bkYellow!40!black] at (112,14) {？};
\end{tikzpicture}
\end{bkfigure}
\end{hatena}

\noindent この章では、「細かく分けて、並べかえたり、たしたりする」という考え方で、いろいろな形の面積と体積を求めます。小学校で学ぶ面積や体積の公式も、すべてこの考え方から作られています。公式を覚えるのではなく、「どう分けて、どう並べかえたか」を理解すれば、わすれても自分で作り直せます。この「細かく分けてたす」考え方は、第5巻の積分（5-9・5-10）につながる、たし合わせの路線（\bkrouteicon[0.9em]{2}）の大切な駅です。

%% ---------------- やってみよう・思い出そう ----------------
\begin{omoidasou}
\bkrecall{1-21}{長方形の面積は「たて×横」。\qty{1}{cm^2} の正方形が何こ並ぶかを数えている。}
\bkrecall{2-7}{体積の単位 \unit{cm^3}・\unit{m^3}。$\qty{1}{L}=\qty{1000}{cm^3}$、$\qty{1}{mL}=\qty{1}{cm^3}$。}
\bkrecall{2-9}{円周＝直径×$\pi$＝$2\pi r$（$r$ は半径）。計算ではふつう $\pi=3.14$ を使う。}
\bkrecall{2-10}{換算は単位つきの1で。\mbox{$\qty{1}{m^2}=\qty{10000}{cm^2}$}、\mbox{$\qty{1}{m^3}=\qty{1000000}{cm^3}$}。}
\bkrecall{4-2}{おうぎ形の弧の長さと面積は、中心角に比例する。}
\bkrecall{4-3}{すいを底面に平行に切ると、断面は底面と相似で、断面の半径（辺）は頂点からの距離に比例する。}
\bkrecall{4-4}{三平方の定理：直角三角形で、斜辺を $c$ とすると $a^2+b^2=c^2$。原点が中心の円は $x^2+y^2=r^2$。}
\end{omoidasou}

\begin{yattemiyou}[円を16に切って並べかえよう]
\noindent\begin{minipage}[c]{0.74\linewidth}
\begin{enumerate}[label={\bkcircled[bkGreen]{\arabic*}}, leftmargin=2.4\zw, itemsep=0pt, topsep=0pt]
\item 画用紙に半径 \qty{5}{cm} の円をかいて、切りぬく。
\item 半分に折ることを4回くり返し、16等分する。
\item 折り目で切り、16まいのおうぎ形にする。
\item とがった所を上・下・上…と交互に並べる。
\item 横の長さと高さをはかり、かけ算しよう。
\end{enumerate}
\end{minipage}\hfill
\begin{minipage}[c]{0.24\linewidth}\centering
\begin{tikzpicture}[x=1mm, y=1mm]
  \bkivfpie{(0,0)}{16}{12}
  \fill[bkInk] (0,0) circle (0.6);
  \draw[bkInk, line width=0.5pt, <->] (0,-14) -- (12,-14);
  \node[font=\sffamily\footnotesize, anchor=north] at (6,-15) {\qty{5}{cm}};
\end{tikzpicture}
\end{minipage}
\end{yattemiyou}

\noindent ソラさんがはかると、横の長さは約 \qty{15.7}{cm}、高さは約 \qty{5}{cm} で、かけ算すると約 \qty{78.5}{cm^2} でした。横の長さ \qty{15.7}{cm} は、何の長さでしょう。円周は「直径×$\pi$」なので $10\times3.14=\qty{31.4}{cm}$（2-9）。横の長さは、ちょうどその半分です。高さの \qty{5}{cm} は、半径と同じです。そのわけは、数学の道具の④で考えます。


%% ---------------- 数学の道具 ----------------
\bksection{doguu}{数学の道具}

\bksubsection{1}{平行四辺形：切って移すと長方形}
\label{bk:4-5:heikou}向かい合う2組の辺が、どちらも平行な四角形を\textbf{平行四辺形}といいます（4-1）。平行四辺形の面積を、長方形に直して求めます。2つの方法で考えましょう。

\bkivfstep{1}{方法1：三角形を切って、反対がわへ移す}
方眼の上に、下の辺が6マス、高さが4マスの平行四辺形をかきます。左はしの直角三角形を切りとって、右はしへ移すと、たて4・横6の長方形になります。

\begin{bkfigure}
\begin{tikzpicture}[x=5mm, y=5mm]
  \draw[bkLine, line width=0.3pt, step=1] (0,0) grid (8,4);
  \fill[bkBlue!20] (2,0) -- (6,0) -- (8,4) -- (2,4) -- cycle;
  \fill[bkOrange!35] (0,0) -- (2,0) -- (2,4) -- cycle;
  \draw[bkInk, line width=1.0pt] (0,0) -- (6,0) -- (8,4) -- (2,4) -- cycle;
  \draw[bkRed, line width=0.9pt, dashed] (2,0) -- (2,4);
  \draw[bk arrow, bkOrange!80!black, line width=0.9pt] (1.2,2.6) to[bend left=25] (6.8,2.6);
  \node[font=\sffamily\footnotesize, anchor=north] at (3,-0.2) {底辺6};
  \node[font=\sffamily\footnotesize, anchor=west, text=bkRed!80!black] at (2.15,1) {高さ4};
  \draw[bk arrow, bkGray, line width=0.9pt] (9,2) -- (11,2);
  \begin{scope}[shift={(12,0)}]
  \draw[bkLine, line width=0.3pt, step=1] (0,0) grid (8,4);
  \fill[bkBlue!20] (2,0) -- (6,0) -- (8,4) -- (2,4) -- cycle;
  \fill[bkOrange!35] (6,0) -- (8,0) -- (8,4) -- cycle;
  \draw[bkInk, line width=1.0pt] (2,0) rectangle (8,4);
  \draw[bkInk!60, line width=0.5pt, dashed] (6,0) -- (8,4);
  \node[font=\sffamily\footnotesize, anchor=north] at (5,-0.2) {横6};
  \node[font=\sffamily\footnotesize, anchor=west] at (8.15,2) {たて4};
  \end{scope}
\end{tikzpicture}
\bkfigcaption{左はしの三角形（オレンジ）を右はしへ移すと、たて4・横6の長方形。面積は $4\times6=24$ マス。}\label{bk:4-5:fig-heikou}
\end{bkfigure}

切って移しただけなので、面積は変わりません。だから、平行四辺形の面積は $6\times4=24$ マスです。マスを数えても、24マスになります（半分のマスは、2つで1マスと数える）。

ここで大事なのは、長方形の「たて」になったのが、ななめの辺ではなく、\textbf{下の辺に垂直にはかった長さ}だということです。平行四辺形で、1つの辺を\textbf{底辺}といい、底辺に垂直にはかった、向かいの辺までの長さを\textbf{高さ}といいます。

\bkivfstep{2}{方法2：細い帯に分けて、ずらす}
\noindent\begin{minipage}[c]{0.56\linewidth}
長方形を、横に細長い帯（はば6、高さ0.5）8本に分けたとします。それぞれの帯を、少しずつ右へずらしていくと、平行四辺形に近い形になります。帯の数をふやして細くするほど、ななめの辺はなめらかになります。

ずらしても、1本1本の帯の面積は変わりません。だから、全部をたした面積も、もとの長方形と同じ $6\times4=24$ です。トランプの山を手でななめにずらしても、カードのまい数が変わらないのと同じです。
\end{minipage}\hfill
\begin{minipage}[c]{0.4\linewidth}\centering
\begin{tikzpicture}[x=5mm, y=5mm]
  \foreach \k in {0,...,7} {\pgfmathsetmacro\xs{\k*0.25}
    \pgfmathparse{int(mod(\k,2))}\ifnum\pgfmathresult=0\def\cc{bkBlue!22}\else\def\cc{bkBlue!40}\fi
    \filldraw[fill=\cc, draw=bkBlue!70!black, line width=0.3pt] (\xs,\k*0.5) rectangle ++(6,0.5);}
  \draw[bkRed, line width=0.8pt, dashed] (0,0) -- (2,4);
  \node[font=\sffamily\footnotesize, anchor=north] at (3,-0.2) {6};
  \node[font=\sffamily\footnotesize, anchor=west] at (8.2,2) {4};
\end{tikzpicture}
\end{minipage}

\vspace{4pt}
2つの方法から、
\bkivfwaku{平行四辺形の面積＝底辺×高さ}
\bkidx{へいこうしへんけいのめんせき}{平行四辺形の面積}

\bksubsection{2}{三角形：同じ三角形を2つ合わせる}
\label{bk:4-5:sankaku}三角形の面積は、同じ三角形をもう1つ用意して、2つを合わせて考えます。

\begin{bkfigure}
\begin{tikzpicture}[x=1mm, y=1mm]
  % 直角三角形2つ → 長方形
  \fill[bkBlue!22] (0,0) -- (30,0) -- (0,20) -- cycle;
  \fill[bkOrange!30] (30,0) -- (30,20) -- (0,20) -- cycle;
  \draw[bkInk, line width=0.9pt] (0,0) rectangle (30,20);
  \draw[bkInk, line width=0.6pt] (30,0) -- (0,20);
  \bkivfright{(0,0)}{(30,0)}{(0,20)}
  \node[font=\sffamily\footnotesize, anchor=north] at (15,-1) {底辺};
  \node[font=\sffamily\footnotesize, anchor=east] at (-1,10) {高さ};
  \node[font=\sffamily\footnotesize, anchor=north, align=center] at (15,-6) {直角三角形2つで\\長方形};
  % ふつうの三角形2つ → 平行四辺形
  \begin{scope}[shift={(48,0)}]
  \fill[bkBlue!22] (0,0) -- (30,0) -- (10,20) -- cycle;
  \fill[bkOrange!30] (30,0) -- (40,20) -- (10,20) -- cycle;
  \draw[bkInk, line width=0.9pt] (0,0) -- (30,0) -- (40,20) -- (10,20) -- cycle;
  \draw[bkInk, line width=0.6pt] (30,0) -- (10,20);
  \draw[bkRed, line width=0.8pt, dashed] (10,0) -- (10,20);
  \bkivfright{(10,0)}{(30,0)}{(10,20)}
  \node[font=\sffamily\footnotesize, anchor=north] at (15,-1) {底辺};
  \node[font=\sffamily\footnotesize, anchor=west, text=bkRed!80!black] at (10.8,6) {高さ};
  \node[font=\sffamily\footnotesize, anchor=north, align=center] at (20,-6) {ふつうの三角形2つで\\平行四辺形};
  \end{scope}
\end{tikzpicture}
\bkfigcaption{青い三角形と同じ三角形（オレンジ）を、180°回して合わせる。}\label{bk:4-5:fig-sankaku}
\end{bkfigure}

\begin{itemize}[leftmargin=1.2\zw, itemsep=2pt, topsep=2pt, label={\color{\bkTT}●}]
\item \textbf{直角三角形}は、2つ合わせると長方形になります。長方形の横は底辺、たては高さです。
\item \textbf{ふつうの三角形}は、2つ合わせると平行四辺形になります。その平行四辺形を、①のように切って移すと長方形になります。平行四辺形の底辺と高さは、もとの三角形の底辺と高さと同じです。
\end{itemize}
どちらも、三角形は長方形（平行四辺形）のちょうど半分です。だから、
\bkivfwaku{三角形の面積＝底辺×高さ÷2}
「÷2」は、同じ三角形を2つ合わせたので、その半分にもどすためです。\bkidx{さんかくけいのめんせき}{三角形の面積}

\bkivfhead{高さが三角形の外に出るとき}
\noindent\begin{minipage}[c]{0.56\linewidth}
角の1つが直角より大きい三角形（鈍角三角形、4-4）では、上の頂点からまっすぐ下ろした線が、底辺の外に出ることがあります。それでも、高さは「底辺をのばした直線までの、垂直な長さ」です。同じ三角形を2つ合わせると、やはり平行四辺形になり、高さも同じなので、面積は「底辺×高さ÷2」で求められます。
\end{minipage}\hfill
\begin{minipage}[c]{0.4\linewidth}\centering
\begin{tikzpicture}[x=1mm, y=1mm]
  \fill[bkBlue!22] (14,0) -- (34,0) -- (2,22) -- cycle;
  \draw[bkInk, line width=0.9pt] (14,0) -- (34,0) -- (2,22) -- cycle;
  \draw[bkGray, line width=0.6pt, dashed] (14,0) -- (0,0);
  \draw[bkRed, line width=0.9pt, dashed] (2,0) -- (2,22);
  \bkivfright{(2,0)}{(14,0)}{(2,22)}
  \node[font=\sffamily\footnotesize, anchor=north] at (24,-1) {底辺};
  \node[font=\sffamily\footnotesize, anchor=east, text=bkRed!80!black] at (0.3,11) {高さ};
\end{tikzpicture}
\end{minipage}

\bkivfhead{面積を変えずに、形を変える（等積変形）}
\noindent\begin{minipage}[c]{0.56\linewidth}
底辺が同じで、高さも同じ三角形は、形がちがっても面積が同じです。上の頂点を、底辺に平行な直線にそって動かしても、高さは変わりません。だから、面積も変わりません。このように、面積を変えずに形を変えることを\textbf{等積変形}といいます。土地の境界線を、面積を変えずにまっすぐに引き直すときなどに使われます。
\end{minipage}\hfill
\begin{minipage}[c]{0.4\linewidth}\centering
\begin{tikzpicture}[x=1mm, y=1mm]
  \draw[bkGray, line width=0.6pt, dashed] (-4,18) -- (46,18);
  \fill[bkBlue!20] (8,0) -- (30,0) -- (2,18) -- cycle;
  \draw[bkBlue!80!black, line width=0.9pt] (8,0) -- (30,0) -- (2,18) -- cycle;
  \fill[bkOrange!25, fill opacity=0.7] (8,0) -- (40,18) -- (30,0) -- cycle;
  \draw[bkOrange!80!black, line width=0.9pt] (8,0) -- (40,18) -- (30,0);
  \fill[bkInk] (2,18) circle (0.8) (40,18) circle (0.8);
  \draw[bk arrow, bkGray, line width=0.7pt] (6,21) -- (36,21);
  \node[font=\sffamily\footnotesize, anchor=north] at (19,-1) {同じ底辺};
\end{tikzpicture}
\end{minipage}

\bksubsection{3}{台形・ひし形・多角形}
\label{bk:4-5:daikei}\bkivfhead{台形：同じ台形を2つ、さかさに合わせる}
\noindent\begin{minipage}[c]{0.54\linewidth}
1組の向かい合う辺が平行な四角形を\textbf{台形}といいます。平行な2辺を上底・下底といい、そのあいだの垂直な長さを高さといいます。

同じ台形を2つ用意し、1つをさかさにして（180°回して）となりに合わせると、平行四辺形になります。この平行四辺形の底辺は「上底＋下底」、高さは台形と同じです。台形はその半分なので、
\end{minipage}\hfill
\begin{minipage}[c]{0.43\linewidth}\centering
\begin{tikzpicture}[x=1mm, y=1mm]
  \fill[bkBlue!22] (0,0) -- (26,0) -- (20,16) -- (6,16) -- cycle;
  \fill[bkOrange!30] (26,0) -- (40,0) -- (46,16) -- (20,16) -- cycle;
  \draw[bkInk, line width=0.9pt] (0,0) -- (40,0) -- (46,16) -- (6,16) -- cycle;
  \draw[bkInk, line width=0.6pt] (26,0) -- (20,16);
  \node[font=\sffamily\footnotesize, anchor=south] at (13,16.5) {上底};
  \node[font=\sffamily\footnotesize, anchor=north] at (13,-0.5) {下底};
  \node[font=\sffamily\footnotesize, anchor=south] at (33,16.5) {下底};
  \node[font=\sffamily\footnotesize, anchor=north] at (33,-0.5) {上底};
  \draw[bkRed, line width=0.8pt, dashed] (12,0) -- (12,16);
  \node[font=\sffamily\footnotesize, anchor=west, text=bkRed!80!black] at (12.6,8) {高さ};
\end{tikzpicture}
\end{minipage}
\bkivfwaku{台形の面積＝（上底＋下底）×高さ÷2}

\bkivfhead{ひし形：長方形の半分}
\noindent\begin{minipage}[c]{0.56\linewidth}
4つの辺の長さがすべて等しい四角形を\textbf{ひし形}といいます。ひし形の2本の対角線は、垂直に交わります。ひし形のまわりに、対角線と平行な辺をもつ長方形をかくと、長方形は8つの直角三角形に分かれ、そのうち4つがひし形です。だから、ひし形は長方形の半分です。
\bkivfline{ひし形の面積＝対角線×対角線÷2}
\end{minipage}\hfill
\begin{minipage}[c]{0.4\linewidth}\centering
\begin{tikzpicture}[x=1mm, y=1mm]
  \draw[bkGray, line width=0.6pt, dashed] (-18,-12) rectangle (18,12);
  \fill[bkBlue!22] (-18,0) -- (0,-12) -- (18,0) -- (0,12) -- cycle;
  \draw[bkInk, line width=0.9pt] (-18,0) -- (0,-12) -- (18,0) -- (0,12) -- cycle;
  \draw[bkInk, line width=0.5pt] (-18,0) -- (18,0) (0,-12) -- (0,12);
  \bkivfright{(0,0)}{(18,0)}{(0,12)}
  \node[font=\sffamily\footnotesize, anchor=north] at (0,-13) {対角線};
  \node[font=\sffamily\footnotesize, anchor=west] at (19.2,0) {対角線};
\end{tikzpicture}
\end{minipage}

\bkivfhead{多角形：三角形に分ける}
\noindent\begin{minipage}[c]{0.56\linewidth}
どんな多角形も、対角線を引けば三角形に分けられます。三角形の面積を1つずつ求めてたせば、全体の面積になります。

たとえば四角形ABCDは、対角線ACで、三角形ABCと三角形ACDに分けられます。どちらも底辺をACにすると、高さはBとDからACまでの垂直な長さです。ACが10、高さが3と5なら、
\bkivfline{$10\times3\div2+10\times5\div2=15+25=40$}
\end{minipage}\hfill
\begin{minipage}[c]{0.4\linewidth}\centering
\begin{tikzpicture}[x=1mm, y=1mm]
  \coordinate (A) at (0,0); \coordinate (C) at (40,0); \coordinate (B) at (14,-12); \coordinate (D) at (26,20);
  \fill[bkBlue!22] (A) -- (B) -- (C) -- cycle;
  \fill[bkOrange!30] (A) -- (C) -- (D) -- cycle;
  \draw[bkInk, line width=0.9pt] (A) -- (B) -- (C) -- (D) -- cycle;
  \draw[bkInk, line width=0.6pt] (A) -- (C);
  \draw[bkRed, line width=0.7pt, dashed] (14,0) -- (B) (26,0) -- (D);
  \node[font=\sffamily\footnotesize, anchor=east] at (-0.8,0) {A};
  \node[font=\sffamily\footnotesize, anchor=north] at (14,-12.8) {B};
  \node[font=\sffamily\footnotesize, anchor=west] at (40.8,0) {C};
  \node[font=\sffamily\footnotesize, anchor=south] at (26,20.8) {D};
  \node[font=\sffamily\footnotesize, anchor=west, text=bkRed!80!black] at (14.6,-6) {3};
  \node[font=\sffamily\footnotesize, anchor=west, text=bkRed!80!black] at (26.6,10) {5};
\end{tikzpicture}
\end{minipage}

\vspace{4pt}
\begin{bkcenter}
\begin{bktable}{colspec={Q[l,m,wd=24mm] Q[l,m,wd=50mm] Q[l,m,wd=44mm]}, rowsep=1.5pt, row{1}={font=\sffamily\small}, cells={font=\small}}
形 & 面積 & どう考えたか \\
長方形 & たて×横 & 1マスが何こ並ぶか（1-21） \\
平行四辺形 & 底辺×高さ & 切って移すと長方形 \\
三角形 & 底辺×高さ÷2 & 2つ合わせると平行四辺形 \\
台形 & （上底＋下底）×高さ÷2 & 2つ合わせると平行四辺形 \\
ひし形 & 対角線×対角線÷2 & 長方形の半分 \\
多角形 & 三角形の和 & 対角線で三角形に分ける \\
\end{bktable}
\bktabcaption{どの公式も、長方形の「たて×横」から作られている。}
\end{bkcenter}

\bksubsection{4}{円の面積：細いおうぎ形に切って並べかえる}
\label{bk:4-5:en}いよいよ円です。はてな？とやってみようのように、円を細いおうぎ形に切って、交互に並べかえます。半径を $r$ とします。

\begin{bkfigure}
\begin{tikzpicture}[x=1mm, y=1mm]
  \bkivfrearr{(0,30)}{8}{12}
  \node[font=\sffamily\footnotesize, anchor=west] at (45,36) {8つに切る};
  \bkivfrearr{(0,13)}{16}{12}
  \node[font=\sffamily\footnotesize, anchor=west] at (45,19) {16に切る};
  \bkivfrearr{(0,-4)}{32}{12}
  \node[font=\sffamily\footnotesize, anchor=west] at (45,2) {32に切る};
  \begin{scope}[shift={(72,-2)}]
  \fill[bkOrange!30] (0,0) rectangle (37.7,12);
  \draw[bkInk, line width=0.9pt] (0,0) rectangle (37.7,12);
  \draw[bkRed, line width=1.2pt] (0,12) -- (37.7,12);
  \node[font=\sffamily\footnotesize, anchor=south, align=center, text=bkRed!80!black] at (18.85,12.8) {円周の半分 $\pi r$};
  \node[font=\sffamily\footnotesize, anchor=west] at (38.5,6) {$r$};
  \node[font=\sffamily\footnotesize, anchor=north] at (18.85,-1) {近づいていく先の長方形};
  \end{scope}
\end{tikzpicture}
\bkfigcaption{切る数をふやすほど、でこぼこが小さくなり、横が $\pi r$・たてが $r$ の長方形に近づく。}\label{bk:4-5:fig-en}
\end{bkfigure}

\bkivfstep{1}{横の長さは、円周の半分}
並べた形の上のでこぼこは、とがった所が下を向いたおうぎ形（オレンジ）の弧です。オレンジは全体の半分なので、上のでこぼこを全部つなぐと、円周の半分になります。下のでこぼこ（青い弧）も円周の半分です。円周は $2\pi r$（2-9）なので、横の長さは $2\pi r\div2=\pi r$ に近づいていきます。

\bkivfstep{2}{たての長さは、半径}
おうぎ形のとがった所から弧までの長さは、どれも半径 $r$ です。細く切るほど、おうぎ形はまっすぐ立つので、形の高さは $r$ に近づきます。

\bkivfstep{3}{だから、円の面積は $\pi r\times r$}
横 $\pi r$、たて $r$ の長方形の面積は、$\pi r\times r=\pi r^2$ です。
\bkivfwaku{円の面積＝半径×半径×$\pi$＝$\pi r^2$}
\bkidx{えんのめんせき}{円の面積}

\bkivfhead{数で確かめる}
半径 \qty{10}{cm} の円で、おうぎ形を「同じ形の二等辺三角形」とみなして並べた形の面積を、コンピュータで計算しました（三角形の高さは、三平方の定理（4-4）で求められます）。
\begin{bkcenter}
\begin{bktable}{colspec={Q[l,m,wd=28mm] *{6}{Q[c,m,wd=12mm]}}, rowsep=1.5pt, cells={font=\small}}
切る数 & 4 & 8 & 16 & 32 & 64 & 128 \\
面積（\unit{cm^2}） & 200.0 & 282.8 & 306.1 & 312.1 & 313.7 & 314.0 \\
\end{bktable}
\bktabcaption{切る数をふやすほど、$\pi\times10^2=314.159\cdots$ に近づいていく。}
\end{bkcenter}
$\pi r^2$ の式では、$3.14159\times10\times10=314.159\cdots$ です。切る数をふやすほど、表の数はこの値に近づいていきます。やってみようの約 \qty{78.5}{cm^2} も、$3.14\times5\times5=78.5$ と合っています。

\bkivfhead{\bkhatenaref の答え}
円を細いおうぎ形に切って並べかえると、細かく切るほど長方形に近づきます。その長方形は、横が円周の半分 $\pi r$、たてが半径 $r$ なので、円の面積は $\pi r^2$ です。どれだけ細かくしても、でこぼこは少しのこるので、ぴったり長方形になるわけではありません。「近づいていく先」の長方形の面積が、円の面積です。

\begin{seikaku}
「どこまでも細かくしたときに、近づいていく先」を考えることを、第5巻では\textbf{極限}とよび、5-3でくわしく学びます。細かく分けたものをたし合わせて、その近づく先を求めるのが、5-10の\textbf{積分}です。円の面積の求め方は、積分の考え方そのものです。
\end{seikaku}

\bkivfhead{おうぎ形の面積}
\noindent\begin{minipage}[c]{0.6\linewidth}
おうぎ形の面積は、中心角に比例します（4-2）。中心角が $a$°のおうぎ形は、円の $\dfrac{a}{360}$ なので、
\bkivfline{おうぎ形の面積＝$\pi r^2\times\dfrac{a}{360}$}
もう1つの見方もあります。おうぎ形を細く切って並べかえると、円と同じように「横が弧の長さの半分、たてが $r$」の長方形に近づきます。だから、
\bkivfline{おうぎ形の面積＝$\dfrac{1}{2}\times$弧の長さ$\times r$}
三角形の「底辺×高さ÷2」で、底辺を弧の長さ、高さを半径にしたのと同じ形です。
\end{minipage}\hfill
\begin{minipage}[c]{0.36\linewidth}\centering
\begin{tikzpicture}[x=1mm, y=1mm]
  \filldraw[fill=bkOrange!30, draw=bkInk, line width=0.9pt] (0,0) -- (20:34) arc (20:80:34) -- cycle;
  \draw[bkRed, line width=1.3pt] (20:34) arc (20:80:34);
  \draw[bkInk, line width=0.5pt] (20:7) arc (20:80:7);
  \node[font=\sffamily\footnotesize] at (50:10.5) {$a$°};
  \node[font=\sffamily\footnotesize, anchor=north west] at (20:17) {$r$};
  \node[font=\sffamily\footnotesize, text=bkRed!80!black, anchor=south west] at (50:34.5) {弧};
\end{tikzpicture}
\end{minipage}

\vspace{3pt}
たとえば、半径 \qty{6}{cm}、中心角60°なら、1つめの式で $\pi\times36\times\dfrac{60}{360}=6\pi$。弧の長さは $2\pi\times6\times\dfrac{60}{360}=2\pi$ なので、2つめの式でも $\dfrac{1}{2}\times2\pi\times6=6\pi$ です。同じになりました。

\bksubsection{5}{柱の体積：底面積×高さ}
\label{bk:4-5:hashira}\bkivfhead{直方体：1辺 \qty{1}{cm} の立方体を数える}
\noindent\begin{minipage}[c]{0.56\linewidth}
たて \qty{3}{cm}・横 \qty{4}{cm}・高さ \qty{2}{cm} の直方体に、1辺 \qty{1}{cm} の立方体（\qty{1}{cm^3}、2-7）がいくつ入るかを数えます。1だんめには、たて3・横4で $3\times4=12$ こ並びます。これが2だん重なるので $12\times2=24$ こ。体積は \qty{24}{cm^3} です。

1だんめの数「$3\times4$」は、底の長方形の面積です。底の面の面積を\textbf{底面積}といいます。
\end{minipage}\hfill
\begin{minipage}[c]{0.4\linewidth}\centering
\begin{tikzpicture}[x={(6mm,0mm)}, y={(3mm,2.2mm)}, z={(0mm,6mm)}]
  \foreach \zz in {0,1} {\foreach \yy in {2,1,0} {\foreach \xx in {0,1,2,3} {\bkivfbox{\xx}{\yy}{\zz}{1}{1}{1}{bkBlue}}}}
  \node[font=\sffamily\footnotesize, anchor=north] at (2,0,-0.15) {横4};
  \node[font=\sffamily\footnotesize, anchor=north west] at (4.1,1.5,0) {たて3};
  \node[font=\sffamily\footnotesize, anchor=east] at (-0.1,0,1) {高さ2};
\end{tikzpicture}
\end{minipage}
\bkivfline{直方体の体積＝たて×横×高さ＝底面積×高さ}

\bkivfhead{角柱・円柱：うすい板を積み重ねる}
\noindent\begin{minipage}[c]{0.56\linewidth}
底面が三角形や円の柱も、同じように考えます。柱を、底面と平行に、とてもうすく輪切りにします。1まい1まいは、底面と同じ形のうすい板で、その体積は「底面積×厚さ」です。柱は、この板を高さの分だけ積み重ねたものなので、
\bkivfwaku{柱の体積＝底面積×高さ}
円柱では、底面積が $\pi r^2$ なので、体積は $\pi r^2h$（$h$ は高さ）です。5-10では、この「うすい円板の体積＝底面積×厚さ」を使って、いろいろな立体の体積を積分で求めます。\bkidx{はしらのたいせき}{柱の体積}
\end{minipage}\hfill
\begin{minipage}[c]{0.4\linewidth}\centering
\begin{tikzpicture}[x=1mm, y=1mm]
  \foreach \k in {0,...,7} {
    \pgfmathparse{int(mod(\k,2))}\ifnum\pgfmathresult=0\def\cc{bkOrange}\else\def\cc{bkYellow}\fi
    \bkivfcyl{(12,\k*3.2)}{10}{3.2}{\cc}}
  \foreach \k in {0,...,7} {\pgfmathsetmacro\xs{\k*1.4}
    \pgfmathparse{int(mod(\k,2))}\ifnum\pgfmathresult=0\def\cc{bkOrange}\else\def\cc{bkYellow}\fi
    \bkivfcyl{(38+\xs,\k*3.2)}{10}{3.2}{\cc}}
  \node[font=\sffamily\footnotesize, anchor=north, align=center] at (12,-4) {コインの山};
  \node[font=\sffamily\footnotesize, anchor=north, align=center] at (43,-4) {ずらした山};
\end{tikzpicture}
\end{minipage}

\vspace{4pt}
コインの山をななめにずらしても、コインのまい数は変わりません。だから、ななめにかたむいた柱の体積も、底面積×高さ（高さは底面に垂直にはかる）です。①の「帯をずらした平行四辺形」の、立体の場合です。

\bksubsection{6}{すいの体積：底面積×高さ×$\dfrac{1}{3}$}
\label{bk:4-5:sui}円すいのように、底面から頂点に向かって細くなっていく立体を\textbf{すい}といいます。底面が四角形なら四角すい、円なら円すいです。すいの体積は、同じ底面・同じ高さの柱の体積の何倍でしょう。3つの方法で調べます。

\bkivfstep{1}{方法1：水を移してはかる}
\noindent\begin{minipage}[c]{0.6\linewidth}
同じ底面・同じ高さの、円すいの容器と円柱の容器を用意します。円すいの容器に水をいっぱい入れて、円柱の容器に移します。これを何回くり返すと、円柱がいっぱいになるでしょう。

やってみると、ちょうど3回でいっぱいになります。円すいの体積は、円柱の $\dfrac{1}{3}$ のようです。
\end{minipage}\hfill
\begin{minipage}[c]{0.36\linewidth}\centering
\begin{tikzpicture}[x=0.8mm, y=0.8mm]
  \bkivfcone{(10,0)}{9}{26}{bkBlue}
  \node[font=\sffamily\footnotesize, anchor=north] at (10,-4) {円すい};
  \bkivfcyl{(40,0)}{9}{26}{bkGray}
  \fill[bkBlue!35] (31,0) arc (180:360:9 and 2.7) -- (49,8.67) arc (0:180:9 and 2.7) -- cycle;
  \filldraw[fill=bkBlue!20, draw=bkBlue!70!black, line width=0.5pt] (40,8.67) ellipse (9 and 2.7);
  \draw[bkGray!60!black, line width=0.6pt] (31,0) -- (31,26) (49,0) -- (49,26) (31,0) arc (180:360:9 and 2.7);
  \draw[bkBlue!70!black, line width=0.4pt, dashed] (31,17.33) -- (49,17.33);
  \node[font=\sffamily\footnotesize, anchor=north] at (40,-4) {円柱};
  \node[font=\sffamily\footnotesize, anchor=west] at (49.5,8.7) {1回め};
\end{tikzpicture}
\end{minipage}

\bkivfstep{2}{方法2：立方体を、3つの同じ四角すいに分ける}
1辺 $a$ の立方体の、1つの頂点Pに注目します。Pをふくまない面は3つあります（下の面・右の面・おくの面）。この3つの面をそれぞれ底面にして、Pを頂点とする四角すいを3つ作ると、立方体がちょうど3つに分かれます。

\begin{bkfigure}
\begin{tikzpicture}[x={(1mm,0mm)}, y={(0.5mm,0.38mm)}, z={(0mm,1mm)}]
  \def\a{16}
  % 立方体と3本の線
  \begin{scope}
  \draw[bkInk!60, line width=0.5pt] (0,0,0) -- (\a,0,0) -- (\a,\a,0) (\a,0,0) -- (\a,0,\a) -- (0,0,\a) -- (0,0,0) (0,0,\a) -- (0,\a,\a) -- (\a,\a,\a) -- (\a,0,\a) (\a,\a,\a) -- (\a,\a,0);
  \draw[bkInk!40, line width=0.5pt, dashed] (0,0,0) -- (0,\a,0) -- (\a,\a,0) (0,\a,0) -- (0,\a,\a);
  \draw[bkRed, line width=0.9pt] (0,0,\a) -- (\a,\a,0);
  \draw[bkRed, line width=0.7pt] (0,0,\a) -- (\a,0,0) (0,0,\a) -- (\a,\a,\a);
  \draw[bkRed, line width=0.7pt, dashed] (0,0,\a) -- (0,\a,0);
  \fill[bkRed] (0,0,\a) circle (0.9);
  \node[font=\sffamily\footnotesize, anchor=south east] at (0,0,\a) {P};
  \node[font=\sffamily\footnotesize, anchor=north] at (9,0,-1) {$a$};
  \end{scope}
  \draw[bk arrow, bkGray, line width=0.8pt] (28,0,11) -- (36,0,11);
  % 3つの四角すい（すべて頂点 P＝(0,0,a)）。うすい点線は、もとの立方体
  \foreach \sx/\col/\nm in {42/bkBlue/下の面が底面, 74/bkOrange/右の面が底面, 106/bkGreen/おくの面が底面} {
  \begin{scope}[shift={(\sx,0,0)}]
  \draw[bkGray!50, line width=0.4pt, dashed] (0,0,0) -- (\a,0,0) -- (\a,\a,0) -- (0,\a,0) -- cycle (0,0,\a) -- (\a,0,\a) -- (\a,\a,\a) -- (0,\a,\a) -- cycle (0,0,0) -- (0,0,\a) (\a,0,0) -- (\a,0,\a) (\a,\a,0) -- (\a,\a,\a) (0,\a,0) -- (0,\a,\a);
  \fill[bkRed] (0,0,\a) circle (0.8);
  \node[font=\sffamily\footnotesize, anchor=south east] at (0,0,\a) {P};
  \node[font=\sffamily\footnotesize, anchor=north] at (8,0,-1) {\nm};
  \end{scope}}
  \begin{scope}[shift={(42,0,0)}]
  \filldraw[fill=bkBlue!35, draw=bkBlue!70!black, line width=0.6pt] (0,0,0) -- (\a,0,0) -- (\a,\a,0) -- (0,\a,0) -- cycle;
  \draw[bkBlue!70!black, line width=0.8pt] (0,0,\a) -- (0,0,0) (0,0,\a) -- (\a,0,0) (0,0,\a) -- (\a,\a,0) (0,0,\a) -- (0,\a,0);
  \end{scope}
  \begin{scope}[shift={(74,0,0)}]
  \filldraw[fill=bkOrange!40, draw=bkOrange!70!black, line width=0.6pt] (\a,0,0) -- (\a,\a,0) -- (\a,\a,\a) -- (\a,0,\a) -- cycle;
  \draw[bkOrange!70!black, line width=0.8pt] (0,0,\a) -- (\a,0,0) (0,0,\a) -- (\a,0,\a) (0,0,\a) -- (\a,\a,\a) (0,0,\a) -- (\a,\a,0);
  \end{scope}
  \begin{scope}[shift={(106,0,0)}]
  \filldraw[fill=bkGreen!30, draw=bkGreen!60!black, line width=0.6pt] (0,\a,0) -- (\a,\a,0) -- (\a,\a,\a) -- (0,\a,\a) -- cycle;
  \draw[bkGreen!60!black, line width=0.8pt] (0,0,\a) -- (0,\a,\a) (0,0,\a) -- (\a,\a,\a) (0,0,\a) -- (0,\a,0) (0,0,\a) -- (\a,\a,0);
  \end{scope}
\end{tikzpicture}
\bkfigcaption{3つの四角すいは、どれも「底面が1辺 $a$ の正方形、高さ $a$」で、合同。3つ合わせると立方体。}\label{bk:4-5:fig-cube3}
\end{bkfigure}

3つの四角すいは、どれも底面が1辺 $a$ の正方形で、頂点Pから底面までの高さも $a$ です。向きがちがうだけで、同じ形・同じ大きさ（合同）です。立方体の体積 $a^3$ を3等分したので、1つの体積は
\bkivfline{$\dfrac{1}{3}a^3=\dfrac{1}{3}\times\underbrace{a\times a}_{\text{底面積}}\times\underbrace{a}_{\text{高さ}}$}
で、ちょうど「底面積×高さ×$\dfrac{1}{3}$」になっています。

\bkivfstep{3}{方法3：うすい板に分けて、たす}
\noindent\begin{minipage}[c]{0.4\linewidth}\centering
\begin{tikzpicture}[x=1mm, y=1mm]
  % 底辺 a=40、高さ h=40 の四角すいを横から見た図。5だんの階段（外側）
  \foreach \k in {1,...,5} {\pgfmathsetmacro\w{40*\k/5}\pgfmathsetmacro\yt{40-8*\k}
    \filldraw[fill=bkOrange!25, draw=bkOrange!70!black, line width=0.4pt] ({20-\w/2},\yt) rectangle ({20+\w/2},\yt+8);}
  \draw[bkBlue!80!black, line width=1.1pt] (0,0) -- (20,40) -- (40,0) -- cycle;
  \node[font=\sffamily\footnotesize, anchor=north] at (20,-1) {$a$};
  \draw[bkGray, line width=0.5pt, <->] (44,0) -- (44,40);
  \node[font=\sffamily\footnotesize, anchor=west] at (44.5,20) {$h$};
\end{tikzpicture}
\end{minipage}\hfill
\begin{minipage}[c]{0.56\linewidth}
底面が1辺 $a$ の正方形、高さ $h$ の四角すいを、底面と平行に $n$ まいのうすい板（厚さ $\dfrac{h}{n}$）に分けます。上から $k$ まいめの板の1辺は、頂点からの距離に比例する（4-3）ので $a\times\dfrac{k}{n}$。その板の体積は
\bkivfline{$\Bigl(a\times\dfrac{k}{n}\Bigr)^2\times\dfrac{h}{n}$}
です。これを $k=1$ から $n$ まで、表計算でたし合わせます。図の階段（$n=5$）は、すいを外側からつつんでいます。
\end{minipage}

\vspace{3pt}
階段の体積が、底面積×高さ（$a^2h$）の何倍になるかを計算しました。板の1辺を1つ小さいほうにずらした「内側の階段」（すいの内側に入る）もならべます。
\begin{bkcenter}
\begin{bktable}{colspec={Q[l,m,wd=40mm] *{4}{Q[c,m,wd=17mm]}}, rowsep=1.5pt, cells={font=\small}}
板の数 $n$ & 5 & 10 & 100 & 1000 \\
外側の階段（$a^2h$ の何倍） & 0.440 & 0.385 & 0.338 & 0.334 \\
内側の階段（$a^2h$ の何倍） & 0.240 & 0.285 & 0.328 & 0.333 \\
\end{bktable}
\bktabcaption{すいの体積は、内側と外側の階段のあいだ。板をふやすと、どちらも $\frac{1}{3}=0.333\cdots$ に近づく。}
\end{bkcenter}
すいは、内側の階段より大きく、外側の階段より小さいので、2つの数のあいだにあります。板をふやすほど、そのはばがせまくなり、0.333…＝$\dfrac{1}{3}$ にはさまれていきます。

3つの方法で、同じ $\dfrac{1}{3}$ が出てきました。方法2は立方体の形のときだけですが、方法1と方法3は、どんな形のすいでも使えます。
\bkivfwaku{すいの体積＝底面積×高さ×$\dfrac{1}{3}$}
\noindent 円すいなら $\dfrac{1}{3}\pi r^2h$、底面が1辺 $a$ の正方形の四角すいなら $\dfrac{1}{3}a^2h$ です。$\dfrac{1}{3}$ は、「すいは先に行くほど細くなるので、柱よりずっと少ない」ことを表しています。5-10では、方法3を積分で計算して、ぴったり $\dfrac{1}{3}$ になることを確かめます。\bkidx{すいのたいせき}{すいの体積}

\bksubsection{7}{球の体積と表面積}
\label{bk:4-5:kyuu}\bkivfhead{半球と「円柱から円すいをくりぬいた形」をくらべる}
半径 $r$ の半球（球の半分）と、底面の半径 $r$・高さ $r$ の円柱から、さかさの円すい（頂点が下の面のまん中）をくりぬいた形をくらべます。どちらも、床からの高さ $y$ の所で、水平にうすく切ってみます。

\begin{bkfigure}
\begin{tikzpicture}[x=2.2mm, y=2.2mm]
  % 左：半球（横から見た断面）r=10、切る高さ y=6
  \fill[bkBlue!18] (-10,0) arc (180:0:10) -- cycle;
  \draw[bkBlue!80!black, line width=0.9pt] (-10,0) arc (180:0:10) -- cycle;
  \draw[bkRed, line width=1.3pt] (-8,6) -- (8,6);
  \draw[bkInk, line width=0.5pt] (0,0) -- (8,6);
  \draw[bkGray, line width=0.5pt, dashed] (0,0) -- (0,6);
  \node[font=\sffamily\footnotesize, anchor=north west] at (3.6,2.7) {$r$};
  \node[font=\sffamily\footnotesize, anchor=east] at (-0.2,3) {$y$};
  \node[font=\sffamily\footnotesize, anchor=south, text=bkRed!80!black] at (-1,6.4) {$\sqrt{r^2-y^2}$};
  \node[font=\sffamily\footnotesize, anchor=north] at (0,-0.6) {半球};
  % 右：円柱からさかさの円すいをくりぬいた形
  \begin{scope}[shift={(28,0)}]
  \fill[bkOrange!25] (-10,0) rectangle (10,10);
  \fill[white] (0,0) -- (-10,10) -- (10,10) -- cycle;
  \draw[bkOrange!80!black, line width=0.9pt] (-10,0) rectangle (10,10);
  \draw[bkOrange!80!black, line width=0.7pt, dashed] (0,0) -- (-10,10) (0,0) -- (10,10);
  \draw[bkRed, line width=1.3pt] (-10,6) -- (-6,6) (6,6) -- (10,6);
  \node[font=\sffamily\footnotesize, anchor=north] at (0,-0.6) {円柱 − 円すい};
  \draw[bkGray, line width=0.4pt, <->] (11.2,0) -- (11.2,10);
  \node[font=\sffamily\footnotesize, anchor=west] at (11.5,5) {$r$};
  \end{scope}
  % 下：切り口を上から見る
  \begin{scope}[shift={(0,-12)}]
  \filldraw[fill=bkBlue!25, draw=bkRed, line width=1pt] (0,0) circle (6.4);
  \node[font=\sffamily\footnotesize, anchor=west, align=left] at (7.2,0) {半径\\$\sqrt{r^2-y^2}$};
  \end{scope}
  \begin{scope}[shift={(28,-12)}]
  \filldraw[fill=bkOrange!30, draw=bkRed, line width=1pt, even odd rule] (0,0) circle (8) (0,0) circle (4.8);
  \node[font=\sffamily\footnotesize, align=center] at (0,0) {あな\\半径 $y$};
  \node[font=\sffamily\footnotesize, anchor=west] at (8.8,0) {外 半径 $r$};
  \end{scope}
\end{tikzpicture}
\bkfigcaption{上は横から見た図、下は高さ $y$ の切り口（赤）を上から見た図。}\label{bk:4-5:fig-hankyu}
\end{bkfigure}

\begin{itemize}[leftmargin=1.2\zw, itemsep=2pt, topsep=2pt, label={\color{\bkTT}●}]
\item \textbf{半球の切り口}は円です。その半径を $x$ とすると、中心から切り口のはしまでが $r$ なので、三平方の定理（4-4）から $x^2+y^2=r^2$、つまり $x^2=r^2-y^2$。切り口の面積は $\pi x^2=\pi(r^2-y^2)$ です。
\item \textbf{円柱−円すいの切り口}は、半径 $r$ の円から、半径 $y$ の円（円すいの切り口）をぬいた輪です。円すいの切り口の半径は、頂点からの高さ $y$ に比例し（4-3）、高さ $r$ で半径 $r$ なので、ちょうど $y$ です。輪の面積は $\pi r^2-\pi y^2=\pi(r^2-y^2)$ です。
\end{itemize}
\textbf{どの高さで切っても、2つの切り口の面積は同じ}です。たとえば $r=10$、$y=6$ なら、半球の切り口は半径 $\sqrt{100-36}=8$ で面積 $64\pi$、輪は $100\pi-36\pi=64\pi$。

2つの立体を、同じ厚さのうすい板に分けると、どの高さの板も体積（切り口の面積×厚さ）が同じです。⑤のコインの山と同じで、板を全部たした体積も同じになります。だから、
\begin{align*}
\text{半球の体積}&=\text{円柱}-\text{円すい}=\pi r^2\times r-\frac{1}{3}\pi r^2\times r&&\text{（高さはどちらも $r$）}\\
&=\pi r^3-\frac{1}{3}\pi r^3=\frac{2}{3}\pi r^3
\end{align*}
球は半球2つ分なので、$\dfrac{2}{3}\pi r^3\times2=\dfrac{4}{3}\pi r^3$ です。
\bkivfwaku{球の体積＝$\dfrac{4}{3}\pi r^3$}
\bkidx{きゅうのたいせき}{球の体積}

\bkivfhead{球の表面積：細いすいに分ける}
\noindent\begin{minipage}[c]{0.6\linewidth}
球の表面を、とても小さな区画に分けます。それぞれの区画と、球の中心を結ぶと、球は、頂点が中心にある、細いすいの集まりになります。どのすいも、高さは半径 $r$ です。

すいの体積は「底面積×高さ×$\dfrac{1}{3}$」なので、全部たすと、
\bkivfline{球の体積＝（底面積の合計）$\times r\times\dfrac{1}{3}$}
底面積の合計は、球の表面積 $S$ です。だから $\dfrac{4}{3}\pi r^3=S\times r\times\dfrac{1}{3}$。両辺に3をかけて $r$ でわると（2-14）、$S=4\pi r^2$ です。
\end{minipage}\hfill
\begin{minipage}[c]{0.36\linewidth}\centering
\begin{tikzpicture}[x=1mm, y=1mm]
  \shade[ball color=bkBlue!30] (0,0) circle (20);
  \fill[bkOrange!45] (0,0) -- (55:20) arc (55:70:20) -- cycle;
  \draw[bkOrange!80!black, line width=0.6pt] (0,0) -- (55:20) arc (55:70:20) -- cycle;
  \fill[bkOrange!45] (0,0) -- (160:20) arc (160:175:20) -- cycle;
  \draw[bkOrange!80!black, line width=0.6pt] (0,0) -- (160:20) arc (160:175:20) -- cycle;
  \fill[bkOrange!45] (0,0) -- (265:20) arc (265:280:20) -- cycle;
  \draw[bkOrange!80!black, line width=0.6pt] (0,0) -- (265:20) arc (265:280:20) -- cycle;
  \fill[bkInk] (0,0) circle (0.7);
  \node[font=\sffamily\footnotesize, anchor=north west] at (1,0) {中心};
  \node[font=\sffamily\footnotesize, anchor=north] at (0,-21.5) {高さ $r$ の細いすい（オレンジ）};
\end{tikzpicture}
\end{minipage}
\bkivfwaku{球の表面積＝$4\pi r^2$}
\noindent 半径 $r$ の円の面積 $\pi r^2$ の、ちょうど4倍です。オレンジの皮を細かくちぎって、同じ半径の円の中にしきつめると、だいたい4つ分になります。家で確かめてみましょう。\bkidx{きゅうのひょうめんせき}{球の表面積}

\bkivfhead{数で確かめる}
半径 \qty{10}{cm} の球の体積は $\dfrac{4}{3}\times3.14\times1000=4186.66\cdots$ で、約 \qty{4190}{cm^3}（約 \qty{4.2}{L}）。表面積は $4\times3.14\times100=\qty{1256}{cm^2}$ です。5-4では、球の体積の式を微分すると、表面積の式になることを確かめます。

\Needspace{11\baselineskip}
\begin{motto}[2]
\textbf{相似な立体の体積は $k^3$ 倍}\par
半径を2倍にすると、円の面積は $\pi(2r)^2=4\pi r^2$ で4倍、球の体積は $\dfrac{4}{3}\pi(2r)^3=8\times\dfrac{4}{3}\pi r^3$ で8倍になります。円どうし・球どうしは相似で、相似比が2なので、面積比は $2^2=4$、体積比は $2^3=8$ です（4-3）。公式の $r^2$、$r^3$ の2と3は、「面積は長さを2回かけたもの、体積は3回かけたもの」（2-7）を表しています。
\end{motto}

\bksubsection{8}{展開図と表面積・回転体}
\label{bk:4-5:tenkai}立体の表面全部の面積を\textbf{表面積}といいます。立体を切り開いて平面に広げた図（\textbf{展開図}）をかくと、表面積は平面の図形の面積の和として求められます。\bkidx{てんかいず}{展開図}

\noindent\begin{minipage}[c]{0.44\linewidth}\centering
\begin{tikzpicture}[x=1mm, y=1mm]
  % 円柱の展開図：r=7（2πr=44）、h=20
  \filldraw[fill=bkBlue!20, draw=bkInk, line width=0.8pt] (0,0) rectangle (44,20);
  \filldraw[fill=bkOrange!30, draw=bkInk, line width=0.8pt] (12,27) circle (7);
  \filldraw[fill=bkOrange!30, draw=bkInk, line width=0.8pt] (12,-7) circle (7);
  \draw[bkInk, line width=0.5pt] (12,27) -- (19,27);
  \node[font=\sffamily\footnotesize, anchor=south] at (15.5,27.8) {$r$};
  \node[font=\sffamily\footnotesize, anchor=north] at (30,-0.8) {$2\pi r$};
  \node[font=\sffamily\footnotesize, anchor=west] at (44.8,10) {$h$};
  \node[font=\sffamily\footnotesize] at (28,10) {側面};
\end{tikzpicture}
\end{minipage}\hfill
\begin{minipage}[c]{0.53\linewidth}
\bkivfhead{円柱}
円柱の展開図は、長方形1つ（側面）と、円2つ（底面）です。側面の長方形の横は、底面の円のまわりに巻きついていたので、円周と同じ $2\pi r$。たては高さ $h$ です。
\bkivfline{表面積＝$2\pi rh+2\times\pi r^2$}
\end{minipage}

\vspace{4pt}
\noindent\begin{minipage}[c]{0.44\linewidth}\centering
\begin{tikzpicture}[x=0.82mm, y=0.82mm]
  % 円すいの展開図：r=3、l=5（l=30、r=18、中心角216°）
  \filldraw[fill=bkBlue!20, draw=bkInk, line width=0.8pt] (0,0) -- (162:30) arc (162:-54:30) -- cycle;
  \draw[bkRed, line width=1.2pt] (162:30) arc (162:-54:30);
  \fill[bkInk] (0,0) circle (0.7);
  \draw[bkInk, line width=0.5pt] (0,0) -- (162:30);
  \node[font=\sffamily\footnotesize, anchor=south] at (162:15) {$\ell$};
  \node[font=\sffamily\footnotesize, anchor=south west, text=bkRed!80!black] at (60:30.5) {弧＝$2\pi r$};
  \filldraw[fill=bkOrange!30, draw=bkInk, line width=0.8pt] ($(-54:30)+(-54:9)$) circle (9);
  \node[font=\sffamily\footnotesize] at ($(-54:30)+(-54:9)$) {底面};
\end{tikzpicture}
\end{minipage}\hfill
\begin{minipage}[c]{0.53\linewidth}
\bkivfhead{円すい}
円すいの展開図は、おうぎ形（側面）と円（底面）です。おうぎ形の半径は、頂点から底面のふちまでの長さ $\ell$（\textbf{母線（ぼせん）}）。おうぎ形の弧は、底面の円周 $2\pi r$ と同じです。④の「おうぎ形の面積＝$\dfrac{1}{2}\times$弧の長さ$\times$半径」から、
\bkivfline{側面積＝$\dfrac{1}{2}\times2\pi r\times\ell=\pi r\ell$}
\end{minipage}

\vspace{4pt}
母線 $\ell$ は、高さ $h$ と底面の半径 $r$ を直角をはさむ辺とする直角三角形の斜辺なので、三平方の定理（4-4）から $\ell=\sqrt{r^2+h^2}$ です。たとえば $r=3$、$h=4$ なら $\ell=\sqrt{9+16}=5$。側面積は $\pi\times3\times5=15\pi$、底面積は $9\pi$ で、表面積は $24\pi$ です。おうぎ形の中心角は、弧が円周 $2\pi\ell$ の $\dfrac{2\pi r}{2\pi\ell}=\dfrac{r}{\ell}$ なので、$360^\circ\times\dfrac{3}{5}=216^\circ$ です。

\bkivfhead{回転体：平面の図形を、軸のまわりに1回転させる}
平面の図形を、1本の直線（軸）のまわりに1回転させてできる立体を\textbf{回転体}といいます。ろくろで作る茶わんや、旋盤でけずる部品は、回転体です。\bkidx{かいてんたい}{回転体}

\begin{bkfigure}
\begin{tikzpicture}[x=1mm, y=1mm]
  % 長方形 → 円柱
  \draw[bkGray, line width=0.6pt, dash dot] (0,-3) -- (0,24);
  \fill[bkBlue!35] (0,0) rectangle (8,18);
  \draw[bkBlue!80!black, line width=0.9pt] (0,0) rectangle (8,18);
  \draw[bk arrow, bkGray, line width=0.7pt] (-3.5,21.5) arc (200:340:3.7 and 1.4);
  \draw[bk arrow, bkGray, line width=0.8pt] (11,9) -- (16,9);
  \bkivfcyl{(27,0)}{8}{18}{bkBlue}
  \node[font=\sffamily\footnotesize, anchor=north] at (17,-5) {長方形 → 円柱};
  % 直角三角形 → 円すい
  \begin{scope}[shift={(46,0)}]
  \draw[bkGray, line width=0.6pt, dash dot] (0,-3) -- (0,24);
  \fill[bkOrange!40] (0,0) -- (8,0) -- (0,18) -- cycle;
  \draw[bkOrange!80!black, line width=0.9pt] (0,0) -- (8,0) -- (0,18) -- cycle;
  \draw[bk arrow, bkGray, line width=0.7pt] (-3.5,21.5) arc (200:340:3.7 and 1.4);
  \draw[bk arrow, bkGray, line width=0.8pt] (11,9) -- (16,9);
  \bkivfcone{(27,0)}{8}{18}{bkOrange}
  \node[font=\sffamily\footnotesize, anchor=north] at (17,-5) {直角三角形 → 円すい};
  \end{scope}
  % 半円 → 球
  \begin{scope}[shift={(92,0)}]
  \draw[bkGray, line width=0.6pt, dash dot] (0,-3) -- (0,24);
  \fill[bkGreen!35] (0,0) arc (-90:90:9) -- cycle;
  \draw[bkGreen!60!black, line width=0.9pt] (0,0) arc (-90:90:9) -- cycle;
  \draw[bk arrow, bkGray, line width=0.7pt] (-3.5,21.5) arc (200:340:3.7 and 1.4);
  \draw[bk arrow, bkGray, line width=0.8pt] (12,9) -- (17,9);
  \shade[ball color=bkGreen!30] (29,9) circle (9);
  \draw[bkGreen!60!black, line width=0.6pt] (29,9) circle (9);
  \draw[bkGreen!60!black, line width=0.4pt] (20,9) arc (180:360:9 and 2.7);
  \draw[bkGreen!60!black, line width=0.4pt, dashed] (20,9) arc (180:0:9 and 2.7);
  \node[font=\sffamily\footnotesize, anchor=north] at (17,-5) {半円 → 球};
  \end{scope}
\end{tikzpicture}
\bkfigcaption{点線が回転の軸。軸にくっついた辺は、立体のまん中の線になる。}\label{bk:4-5:fig-kaiten}
\end{bkfigure}

軸からいちばん遠い所が、回転体の半径になります。直角三角形を、直角をはさむ1辺を軸にして回すと、その辺が円すいの高さ、もう1つの辺が底面の半径になります。半円を、直径を軸にして回すと球です。

\bksubsection{9}{方眼で数えて、細かくするほど正確に}
\label{bk:4-5:komakaku}円の面積は、別の方法でも確かめられます。半径 \qty{10}{cm} の円を、1マス \qty{1}{cm} の方眼にかいて、マスを数えます。

\noindent\begin{minipage}[c]{0.44\linewidth}\centering
\begin{tikzpicture}[x=4mm, y=4mm]
  \foreach \i/\j in {0/0,0/1,0/2,0/3,0/4,0/5,0/6,0/7,0/8,1/0,1/1,1/2,1/3,1/4,1/5,1/6,1/7,1/8,2/0,2/1,2/2,2/3,2/4,2/5,2/6,2/7,2/8,3/0,3/1,3/2,3/3,3/4,3/5,3/6,3/7,3/8,4/0,4/1,4/2,4/3,4/4,4/5,4/6,4/7,5/0,5/1,5/2,5/3,5/4,5/5,5/6,5/7,6/0,6/1,6/2,6/3,6/4,6/5,6/6,7/0,7/1,7/2,7/3,7/4,7/5,8/0,8/1,8/2,8/3}
    {\fill[bkBlue!30] (\i,\j) rectangle ++(1,1);}
  \foreach \i/\j in {0/9,1/9,2/9,3/9,4/8,4/9,5/8,6/7,7/6,7/7,8/4,8/5,9/0,9/1,9/2,9/3,9/4}
    {\fill[bkOrange!40] (\i,\j) rectangle ++(1,1);}
  \draw[bkLine, line width=0.3pt, step=1] (0,0) grid (10,10);
  \draw[bkRed, line width=1.2pt] (10,0) arc (0:90:10);
  \draw[bkInk, line width=0.6pt] (0,0) -- (10,0) (0,0) -- (0,10);
  \node[font=\sffamily\footnotesize, anchor=north] at (5,-0.2) {\qty{10}{cm}};
\end{tikzpicture}
\end{minipage}\hfill
\begin{minipage}[c]{0.53\linewidth}
図は、円の $\dfrac{1}{4}$ です。
\begin{itemize}[leftmargin=1.2\zw, itemsep=1pt, topsep=2pt, label={\color{\bkTT}●}]
\item 全部が円の内側にあるマス（青）：69こ
\item 円のふちにかかるマス（オレンジ）：17こ
\end{itemize}
円全体ではこの4倍です。だから、円の面積は
\bkivfline{青だけの $69\times4=\qty{276}{cm^2}$ より大きく、}
\bkivfline{青とオレンジの $(69+17)\times4=\qty{344}{cm^2}$ より小さい}
とわかります。
\end{minipage}

\vspace{4pt}
マスを小さくすると、ふちにかかるマスの面積がへって、はさむはばがせまくなります。コンピュータで数えた結果です。
\begin{bkcenter}
\begin{bktable}{colspec={Q[l,m,wd=40mm] *{4}{Q[c,m,wd=17mm]}}, rowsep=1.5pt, cells={font=\small}}
1マスの1辺（cm） & 1 & 0.5 & 0.25 & 0.1 \\
内側のマスだけ（\unit{cm^2}） & 276 & 294 & 304 & 310.2 \\
ふちまでふくむ（\unit{cm^2}） & 344 & 331 & 323.3 & 318.0 \\
\end{bktable}
\bktabcaption{マスを細かくするほど、2つの数は $\pi r^2=314.159\cdots$ をはさんで近づく。}
\end{bkcenter}

④の「並べかえ」と、⑨の「方眼で数える」は、まったくちがう方法ですが、どちらも同じ $\pi r^2$ に近づきます。⑥の方法3でも、「外側の階段」と「内側の階段」で、すいの体積をはさみました。

\Needspace{16\baselineskip}
\begin{motto}[2]
\textbf{切り口の面積が同じなら、体積も同じ}\par
\noindent\begin{minipage}[c]{0.6\linewidth}
⑤のコインの山と、⑦の半球で使った考えをまとめます。2つの立体を、同じ高さのところで水平に切ります。\textbf{どの高さで切っても切り口の面積が同じなら、2つの立体の体積は同じ}です。うすい板に分けると、どの高さの板も体積が同じだからです。400年ほど前のイタリアの数学者の名前から、これを\textbf{カヴァリエリの原理}といいます。

たとえば、底面積と高さが同じ四角すいと円すいは、体積が同じです。どちらも、切り口の面積は「頂点からの距離の2乗」に比例して、同じようにふえていくからです（4-3）。
\end{minipage}\hfill
\begin{minipage}[c]{0.36\linewidth}\centering
\begin{tikzpicture}[x=0.72mm, y=0.72mm]
  % 四角すい（横から）と円すい（横から）。底面積が同じ（正方形の1辺 a、円の半径 r で a²＝πr²）
  \filldraw[fill=bkBlue!18, draw=bkBlue!70!black, line width=0.7pt] (0,0) -- (18,0) -- (24,5) -- (12,28) -- cycle;
  \draw[bkBlue!70!black, line width=0.6pt] (18,0) -- (12,28);
  \draw[bkBlue!70!black, line width=0.4pt, dashed] (0,0) -- (6,5) -- (24,5) (6,5) -- (12,28);
  \bkivfcone{(42,0)}{11.3}{28}{bkOrange}
  \draw[bkRed, line width=1pt, dashed] (-3,14) -- (56,14);
  \node[font=\sffamily\footnotesize, anchor=north] at (12,-2) {四角すい};
  \node[font=\sffamily\footnotesize, anchor=north] at (42,-4.5) {円すい};
  \node[font=\sffamily\footnotesize, anchor=south, text=bkRed!80!black] at (27,30) {同じ高さで切る};
\end{tikzpicture}
\end{minipage}
\end{motto}

\begin{taisetsu}
細かく分けて、たし合わせる。分け方を細かくするほど、答えは正確になる。
\end{taisetsu}

\noindent この考え方は、形がふくざつで公式がない図形（池の面積、葉の面積、臓器の体積など）にも使えます。5-9では、グラフの下の面積を細い長方形の和で求め（区分求積）、5-10で、その近づく先を計算する方法（積分）を学びます。

\begin{taisetsu}
\def\tsl#1{\makebox[5.2\zw][l]{\sffamily\bfseries\small\color{bkDeep}#1}}
\tsl{平行四辺形}底辺×高さ\qquad\tsl{三角形}底辺×高さ÷2\\[3pt]
\tsl{台形}（上底＋下底）×高さ÷2\qquad\tsl{ひし形}対角線×対角線÷2\\[3pt]
\tsl{円}$\pi r^2$\qquad\tsl{おうぎ形}$\dfrac{1}{2}\times$弧の長さ$\times r$\\[3pt]
\tsl{柱}底面積×高さ\qquad\tsl{すい}底面積×高さ×$\dfrac{1}{3}$\\[3pt]
\tsl{球}体積 $\dfrac{4}{3}\pi r^3$\qquad 表面積 $4\pi r^2$
\end{taisetsu}

\begin{kotoba}
\small
\bkword{底辺・高さ}{1つの辺を底辺とするとき、底辺（をのばした直線）に垂直にはかった、向かいの頂点や辺までの長さが高さ。}
\bkword{等積変形}{面積を変えずに形を変えること。}
\bkword{底面積}{柱やすいの底の面の面積。}
\bkword{すい}{底面から1つの頂点に向かって細くなる立体。四角すい・円すいなど。}
\bkword{母線}{円すいの頂点と、底面の円のふちを結ぶ線分。}
\bkword{表面積・展開図}{立体の表面全部の面積。立体を切り開いて平面に広げた図。}
\bkword{回転体}{平面の図形を、軸のまわりに1回転させてできる立体。}
\end{kotoba}

\bkkigocards{\bkkigocard{$S$・$V$}{エス・ブイ}{面積を $S$、体積を $V$ で表すことが多い。物理では面積に $A$ も使う}{surface, volume}{4-5}}{}

\begin{column}{アルキメデスのお墓}
2200年ほど前のギリシャの数学者アルキメデスは、「球の体積は、それがぴったり入る円柱の体積の $\dfrac{2}{3}$、球の表面積も円柱の表面積の $\dfrac{2}{3}$」であることを見つけました。円柱は半径 $r$・高さ $2r$ なので体積は $2\pi r^3$、球の体積 $\dfrac{4}{3}\pi r^3$ はその $\dfrac{2}{3}$ です。アルキメデスはこの発見がとても気に入っていて、お墓に「円柱の中に入った球」の図をきざんでほしいとたのんだと伝えられています。
\end{column}

%% ---------------- 例題 ----------------
\Needspace{28\baselineskip}
\bksection{reidai}{例題でたしかめよう}
\noindent 例題では、$\pi$ をふくむ答えは「$36\pi$\,cm$^2$」のように $\pi$ のまま書き、およその大きさを $\pi=3.14$ として「約」をつけて添えます。

\begin{reidai}[平行四辺形と三角形の面積]
次の面積を求めましょう。\par
（1）底辺 \qty{8}{cm}、高さ \qty{5}{cm} の平行四辺形\quad（2）底辺 \qty{6}{cm}、高さ \qty{7}{cm} の三角形\par
（3）底辺 \qty{4}{cm}、高さ \qty{6}{cm} の三角形（高さが三角形の外に出る）\par
\kangaekata 平行四辺形は底辺×高さ（①）。三角形はその半分なので÷2（②）。高さが外に出ても、計算は同じ。\par
\shiki（1）$8\times5=40$\par
\hspace*{2.6\zw}（2）$6\times7\div2=42\div2=21$\par
\hspace*{2.6\zw}（3）$4\times6\div2=24\div2=12$\par
\kotae（1）\qty{40}{cm^2}\quad（2）\qty{21}{cm^2}\quad（3）\qty{12}{cm^2}
\end{reidai}

\begin{reidai}[台形とひし形の面積]
（1）上底 \qty{4}{cm}、下底 \qty{10}{cm}、高さ \qty{6}{cm} の台形と、（2）対角線が \qty{8}{cm} と \qty{6}{cm} のひし形の面積を求めましょう。\par
\kangaekata 台形は（上底＋下底）×高さ÷2、ひし形は対角線×対角線÷2（③）。\par
\shiki（1）$(4+10)\times6\div2=14\times6\div2=42$\par
\hspace*{2.6\zw}（2）$8\times6\div2=24$\par
\kotae（1）\qty{42}{cm^2}\quad（2）\qty{24}{cm^2}
\end{reidai}

\noindent\begin{minipage}[t]{0.485\linewidth}\vspace{0pt}\raggedright
\begin{jibun}\raggedright
次の面積を求めましょう。\par（1）底辺 \qty{12}{cm}、高さ \qty{7}{cm} の平行四辺形\par（2）底辺 \qty{9}{cm}、高さ \qty{4}{cm} の三角形\par（3）底辺 \qty{5}{cm}、高さ \qty{8}{cm} の三角形（高さが外に出る）
\bkjans{（1）\nobreak \qty{84}{cm^2}。（2）\nobreak \qty{18}{cm^2}。（3）\nobreak \qty{20}{cm^2}}
\end{jibun}
\end{minipage}\hfill
\begin{minipage}[t]{0.485\linewidth}\vspace{0pt}\raggedright
\begin{jibun}\raggedright
（1）上底 \qty{3}{cm}、下底 \qty{7}{cm}、高さ \qty{5}{cm} の台形と、（2）対角線が \qty{10}{cm} と \qty{7}{cm} のひし形の面積を求めましょう。
\bkjans{（1）\nobreak $(3+7)\times5\div2=\qty{25}{cm^2}$。（2）\nobreak $10\times7\div2=\qty{35}{cm^2}$}
\end{jibun}
\end{minipage}

\begin{reidai}[四角形を三角形に分ける]
四角形ABCDの対角線ACの長さは \qty{10}{cm} です。BとDから対角線ACまでの垂直な長さは、それぞれ \qty{3}{cm} と \qty{5}{cm} です。四角形ABCDの面積を求めましょう。\par
\kangaekata 対角線ACで、三角形ABCと三角形ACDに分ける。どちらも底辺はAC（③）。\par
\shiki 三角形ABC：$10\times3\div2=15$\par
\hspace*{2.6\zw}三角形ACD：$10\times5\div2=25$\par
\hspace*{2.6\zw}合わせて $15+25=40$\par
\kotae \qty{40}{cm^2}
\end{reidai}

\begin{reidai}[円とおうぎ形の面積]
（1）半径 \qty{6}{cm} の円の面積と、（2）半径 \qty{6}{cm}、中心角60°のおうぎ形の面積を求めましょう。\par
\kangaekata 円は $\pi r^2$（④）。おうぎ形は、円の $\dfrac{\text{中心角}}{360}$。\par
\shiki（1）$\pi\times6^2=36\pi$\par
\hspace*{4.6\zw}$36\times3.14=113.04$\par
\hspace*{2.6\zw}（2）$36\pi\times\dfrac{60}{360}=36\pi\times\dfrac{1}{6}=6\pi$\par
\hspace*{4.6\zw}$6\times3.14=18.84$\par
\kotae（1）$36\pi$\,cm$^2$（約 \qty{113}{cm^2}）\quad（2）$6\pi$\,cm$^2$（約 \qty{18.8}{cm^2}）
\end{reidai}

\noindent\begin{minipage}[t]{0.485\linewidth}\vspace{0pt}\raggedright
\begin{jibun}\raggedright
四角形ABCDの対角線ACは \qty{12}{cm}、BとDからACまでの垂直な長さは \qty{4}{cm} と \qty{2}{cm} です。面積を求めましょう。
\bkjans{$12\times4\div2+12\times2\div2=24+12=\qty{36}{cm^2}$}
\end{jibun}
\end{minipage}\hfill
\begin{minipage}[t]{0.485\linewidth}\vspace{0pt}\raggedright
\begin{jibun}\raggedright
（1）半径 \qty{5}{cm} の円の面積と、（2）半径 \qty{10}{cm}、中心角72°のおうぎ形の面積を求めましょう。
\bkjans{（1）\nobreak $25\pi$\,cm$^2$（約 \qty{78.5}{cm^2}）。（2）\nobreak $100\pi\times\dfrac{72}{360}=20\pi$\,cm$^2$（約 \qty{62.8}{cm^2}）}
\end{jibun}
\end{minipage}

\bkivfhead{$\pi$ のまま答えるわけ}
\noindent\begin{minipage}[c]{0.62\linewidth}
$\pi$ は $3.14159\cdots$ と、どこまでも続く数です（2-9）。3.14 で計算すると、答えに小さなずれが入ります。$36\pi$ のように $\pi$ のまま書いておけば、ずれのない正確な答えになり、あとで別の計算に使うときも便利です。

中心角60°のおうぎ形は、右の図のように、円をピザのように6等分した1切れです（$360\div60=6$）。だから面積も円の $\dfrac{1}{6}$ で、$36\pi\div6=6\pi$ です。中心角が何分の1かを考えると、計算の見通しが立ちます。
\end{minipage}\hfill
\begin{minipage}[c]{0.34\linewidth}\centering
\begin{tikzpicture}[x=1mm, y=1mm]
  \foreach \k in {0,...,5} {\filldraw[fill=bkBlue!15, draw=bkInk!60, line width=0.4pt] (0,0) -- (\k*60:16) arc (\k*60:\k*60+60:16) -- cycle;}
  \filldraw[fill=bkOrange!40, draw=bkInk, line width=0.8pt] (0,0) -- (0:16) arc (0:60:16) -- cycle;
  \draw[bkInk, line width=0.5pt] (0:5) arc (0:60:5);
  \node[font=\sffamily\footnotesize] at (30:8.5) {60°};
  \node[font=\sffamily\footnotesize, anchor=north] at (0,-17.5) {円の $\dfrac{1}{6}$};
\end{tikzpicture}
\end{minipage}

\begin{reidai}[柱の体積]
（1）底面が、直角をはさむ辺 \qty{3}{cm} と \qty{4}{cm} の直角三角形で、高さ \qty{10}{cm} の三角柱と、（2）底面の半径 \qty{3}{cm}、高さ \qty{10}{cm} の円柱の体積を求めましょう。\par
\kangaekata 柱の体積＝底面積×高さ（⑤）。まず底面積を求める。\par
\shiki（1）底面積 $3\times4\div2=6$、体積 $6\times10=60$\par
\hspace*{2.6\zw}（2）底面積 $\pi\times3^2=9\pi$、体積 $9\pi\times10=90\pi$\par
\hspace*{4.6\zw}$90\times3.14=282.6$\par
\kotae（1）\qty{60}{cm^3}\quad（2）$90\pi$\,cm$^3$（約 \qty{283}{cm^3}）
\end{reidai}

\begin{reidai}[すいの体積]
（1）底面が1辺 \qty{6}{cm} の正方形で、高さ \qty{10}{cm} の四角すいと、（2）底面の半径 \qty{3}{cm}、高さ \qty{4}{cm} の円すいの体積を求めましょう。\par
\kangaekata すいの体積＝底面積×高さ×$\dfrac{1}{3}$（⑥）。$\dfrac{1}{3}$ をわすれない。\par
\shiki（1）$6\times6\times10\times\dfrac{1}{3}=360\times\dfrac{1}{3}=120$\par
\hspace*{2.6\zw}（2）$\pi\times3^2\times4\times\dfrac{1}{3}=36\pi\times\dfrac{1}{3}=12\pi$\par
\hspace*{4.6\zw}$12\times3.14=37.68$\par
\kotae（1）\qty{120}{cm^3}\quad（2）$12\pi$\,cm$^3$（約 \qty{37.7}{cm^3}）
\end{reidai}

\noindent\begin{minipage}[t]{0.485\linewidth}\vspace{0pt}\raggedright
\begin{jibun}\raggedright
（1）底面が、底辺 \qty{6}{cm}・高さ \qty{5}{cm} の三角形で、高さ \qty{8}{cm} の三角柱と、（2）底面の半径 \qty{2}{cm}、高さ \qty{5}{cm} の円柱の体積を求めましょう。
\bkjans{（1）\nobreak $15\times8=\qty{120}{cm^3}$。（2）\nobreak $4\pi\times5=20\pi$\,cm$^3$（約 \qty{62.8}{cm^3}）}
\end{jibun}
\end{minipage}\hfill
\begin{minipage}[t]{0.485\linewidth}\vspace{0pt}\raggedright
\begin{jibun}\raggedright
（1）底面が1辺 \qty{5}{cm} の正方形で、高さ \qty{6}{cm} の四角すいと、（2）底面の半径 \qty{6}{cm}、高さ \qty{5}{cm} の円すいの体積を求めましょう。
\bkjans{（1）\nobreak $25\times6\times\dfrac{1}{3}=\qty{50}{cm^3}$。（2）\nobreak $36\pi\times5\times\dfrac{1}{3}=60\pi$\,cm$^3$（約 \qty{188}{cm^3}）}
\end{jibun}
\end{minipage}

\begin{reidai}[球の体積と表面積]
半径 \qty{6}{cm} の球の体積と表面積を求めましょう。\par
\kangaekata 体積 $\dfrac{4}{3}\pi r^3$、表面積 $4\pi r^2$（⑦）。$r^3$ は $r$ を3回かける。\par
\shiki 体積 $\dfrac{4}{3}\times\pi\times6^3=\dfrac{4}{3}\times216\pi=288\pi$\par
\hspace*{4.6\zw}$288\times3.14=904.32$\par
\hspace*{2.6\zw}表面積 $4\times\pi\times6^2=4\times36\pi=144\pi$\par
\hspace*{4.6\zw}$144\times3.14=452.16$\par
\kotae 体積 $288\pi$\,cm$^3$（約 \qty{904}{cm^3}）、表面積 $144\pi$\,cm$^2$（約 \qty{452}{cm^2}）
\end{reidai}

\bkivfhead{表面積は、同じ半径の円4つ分}
\noindent\begin{minipage}[c]{0.56\linewidth}
例題7の表面積 $144\pi$ は、半径 \qty{6}{cm} の円の面積 $36\pi$ のちょうど4倍です（$4\times36\pi=144\pi$）。⑦のオレンジの皮の実験のように、球の皮をはがして平らに広げると、同じ半径の円4つ分になります。

体積の答えは \unit{cm^3}、表面積の答えは \unit{cm^2} です。単位の3と2は、長さを3回かけたか、2回かけたかを表しています（2-7）。
\end{minipage}\hfill
\begin{minipage}[c]{0.4\linewidth}\centering
\begin{tikzpicture}[x=0.85mm, y=0.85mm]
  \shade[ball color=bkOrange!50] (0,0) circle (9);
  \node[font=\sffamily\footnotesize, anchor=north] at (0,-10.5) {球};
  \node[font=\Large] at (13,0) {$=$};
  \foreach \xx/\yy in {25/5, 41/5, 25/-11, 41/-11} {\filldraw[fill=bkOrange!30, draw=bkOrange!70!black, line width=0.6pt] (\xx,\yy+3) circle (7);}
  \node[font=\sffamily\footnotesize, anchor=north] at (33,-16.5) {同じ半径の円4つ};
\end{tikzpicture}
\end{minipage}

\begin{reidai}[円柱と円すいの表面積]
（1）底面の半径 \qty{3}{cm}、高さ \qty{5}{cm} の円柱と、（2）底面の半径 \qty{3}{cm}、高さ \qty{4}{cm} の円すいの表面積を求めましょう。\par
\kangaekata 展開図で考える（⑧）。円柱の側面は横 $2\pi r$・たて $h$ の長方形。円すいの側面は $\pi r\ell$ で、母線 $\ell$ は三平方の定理で求める。\par
\shiki（1）側面 $2\pi\times3\times5=30\pi$、底面2つ $2\times9\pi=18\pi$\par
\hspace*{4.6\zw}$30\pi+18\pi=48\pi$\par
\hspace*{2.6\zw}（2）母線 $\ell=\sqrt{3^2+4^2}=\sqrt{25}=5$\par
\hspace*{4.6\zw}側面 $\pi\times3\times5=15\pi$、底面 $9\pi$\par
\hspace*{4.6\zw}$15\pi+9\pi=24\pi$\par
\kotae（1）$48\pi$\,cm$^2$（約 \qty{151}{cm^2}）\quad（2）$24\pi$\,cm$^2$（約 \qty{75.4}{cm^2}）
\end{reidai}

\noindent\begin{minipage}[t]{0.485\linewidth}\vspace{0pt}\raggedright
\begin{jibun}\raggedright
半径 \qty{10}{cm} の球の体積と表面積を求めましょう。
\bkjans{体積 $\dfrac{4}{3}\pi\times1000=\dfrac{4000}{3}\pi$\,cm$^3$（約 \qty{4190}{cm^3}）。表面積 $400\pi$\,cm$^2$（約 \qty{1256}{cm^2}）}
\end{jibun}
\end{minipage}\hfill
\begin{minipage}[t]{0.485\linewidth}\vspace{0pt}\raggedright
\begin{jibun}\raggedright
（1）底面の半径 \qty{2}{cm}、高さ \qty{6}{cm} の円柱と、（2）底面の半径 \qty{5}{cm}、高さ \qty{12}{cm} の円すいの表面積を求めましょう。
\bkjans{（1）\nobreak $24\pi+8\pi=32\pi$\,cm$^2$（約 \qty{100}{cm^2}）。（2）\nobreak $\ell=\sqrt{25+144}=13$、$65\pi+25\pi=90\pi$\,cm$^2$（約 \qty{283}{cm^2}）}
\end{jibun}
\end{minipage}

\begin{reidai}[回転体の体積]
（1）たて \qty{4}{cm}、横 \qty{3}{cm} の長方形を、たての辺を軸にして1回転させた立体と、（2）直角をはさむ辺が \qty{6}{cm} と \qty{8}{cm} の直角三角形を、\qty{8}{cm} の辺を軸にして1回転させた立体の体積を求めましょう。\par
\kangaekata どんな立体になるか、軸から遠い所が半径になることを、図をかいて確かめる（⑧）。\par
\shiki（1）半径3・高さ4の円柱。$\pi\times3^2\times4=36\pi$\par
\hspace*{2.6\zw}（2）半径6・高さ8の円すい。$\pi\times6^2\times8\times\dfrac{1}{3}=96\pi$\par
\kotae（1）$36\pi$\,cm$^3$（約 \qty{113}{cm^3}）\quad（2）$96\pi$\,cm$^3$（約 \qty{301}{cm^3}）
\end{reidai}

\begin{reidai}[タンクや水そうの容量を L で]
（1）直径 \qty{2}{m}、高さ \qty{1.5}{m} の円柱形のタンクに入る水は何Lですか（$\pi=3.14$）。\par
（2）内側の長さが、たて \qty{30}{cm}・横 \qty{60}{cm}・深さ \qty{36}{cm} の水そうに入る水は何Lですか。\par
\kangaekata 半径は直径の半分。体積を求めてから、単位つきの1で換算する（2-10）。$\qty{1}{m^3}=100\times100\times\qty{100}{cm^3}=\qty{1000000}{cm^3}$、$\qty{1}{L}=\qty{1000}{cm^3}$ なので、$\qty{1}{m^3}=\qty{1000}{L}$。\par
\shiki（1）半径1。$\pi\times1^2\times1.5=1.5\pi=4.71$（\unit{m^3}）\par
\hspace*{4.6\zw}$\qty{4.71}{m^3}\times\qty{1000}{L/m^3}=\qty{4710}{L}$\par
\hspace*{2.6\zw}（2）$30\times60\times36=\num{64800}$（\unit{cm^3}）\par
\hspace*{4.6\zw}$\qty{64800}{cm^3}\div\qty{1000}{cm^3/L}=\qty{64.8}{L}$\par
\kotae（1）約 \qty{4710}{L}\quad（2）\qty{64.8}{L}
\end{reidai}

\noindent\begin{minipage}[t]{0.485\linewidth}\vspace{0pt}\raggedright
\begin{jibun}\raggedright
（1）たて \qty{5}{cm}、横 \qty{2}{cm} の長方形を、たての辺を軸にして1回転させた立体と、（2）直角をはさむ辺が \qty{3}{cm} と \qty{4}{cm} の直角三角形を、\qty{4}{cm} の辺を軸にして1回転させた立体の体積を求めましょう。
\bkjans{（1）\nobreak 半径2・高さ5の円柱で $20\pi$\,cm$^3$（約 \qty{62.8}{cm^3}）。（2）\nobreak 半径3・高さ4の円すいで $12\pi$\,cm$^3$（約 \qty{37.7}{cm^3}）}
\end{jibun}
\end{minipage}\hfill
\begin{minipage}[t]{0.485\linewidth}\vspace{0pt}\raggedright
\begin{jibun}\raggedright
（1）直径 \qty{1}{m}、高さ \qty{2}{m} の円柱形のタンクに入る水は何Lですか（$\pi=3.14$）。\par（2）内側が、たて \qty{25}{cm}・横 \qty{50}{cm}・深さ \qty{30}{cm} の水そうに入る水は何Lですか。
\bkjans{（1）\nobreak $\pi\times0.5^2\times2=0.5\pi=\qty{1.57}{m^3}$ で、約 \qty{1570}{L}。（2）\nobreak \qty{37500}{cm^3} で \qty{37.5}{L}}
\end{jibun}
\end{minipage}

\bksection{tsumazuki}{まちがえやすいところをチェック}

\begin{tsumazuki}{ななめの辺を高さにする}
底辺 \qty{8}{cm}、ななめの辺 \qty{6}{cm}、高さ \qty{5}{cm} の平行四辺形の面積は？\par
\bkng{$8\times6=\qty{48}{cm^2}$}
\bkok{$8\times5=\qty{40}{cm^2}$}
高さは、底辺に垂直にはかった長さです。ななめの辺は、高さより長くなります。①の図で、長方形のたてになったのは高さでした。
\end{tsumazuki}

\begin{tsumazuki}{÷2や $\dfrac{1}{3}$ をわすれる}
底面の半径 \qty{3}{cm}、高さ \qty{4}{cm} の円すいの体積は？\par
\bkng{$\pi\times3^2\times4=36\pi$}
\bkok{$\pi\times3^2\times4\times\dfrac{1}{3}=12\pi$}
$36\pi$ は、同じ底面・同じ高さの円柱の体積です。円すいは、とがっている分、ずっと小さくなります。三角形の÷2も同じで、「2つ合わせると平行四辺形」を思い出しましょう。
\end{tsumazuki}

\begin{tsumazuki}{直径を半径として使う}
直径 \qty{10}{cm} の円の面積は？\par
\bkng{$\pi\times10^2=100\pi$}
\bkok{半径は5なので、$\pi\times5^2=25\pi$}
$\pi r^2$ の $r$ は半径です。問題に「直径」とあったら、まず半分にして半径を求めましょう。直径のまま計算すると、4倍の大きさになってしまいます。
\end{tsumazuki}

\begin{tsumazuki}{円すいの表面積に、高さを使う}
\bkng{側面積＝$\pi\times\text{半径}\times\text{高さ}$}
\bkok{側面積＝$\pi\times\text{半径}\times\text{母線}$}
側面のおうぎ形の半径は、頂点から底面のふちまでの母線 $\ell$ です。高さ $h$ ではありません。$\ell=\sqrt{r^2+h^2}$ で、母線のほうが高さより長くなります。
\end{tsumazuki}

\begin{tsumazuki}{$\unit{m^3}$ と L の換算}
\bkng{$\qty{1}{m^3}=\qty{100}{L}$}
\bkok{$\qty{1}{m^3}=\qty{1000}{L}$}
\qty{1}{m}＝\qty{100}{cm} なので、100を3回かけて \mbox{$\qty{1}{m^3}=\qty{1000000}{cm^3}$}。これを \qty{1}{L} の \qty{1000}{cm^3} でわると、1000です。体積では、長さの換算の数を3回かけます（2-10）。
\end{tsumazuki}


%% ---------------- 物理・工学 ----------------
\begin{butsuri}
\par\addvspace{3pt}\textbf{パイプの断面積と流量}\par\noindent
\begin{minipage}[c]{0.62\linewidth}
水道管の中を、水が1秒に \qty{2}{m} の速さで流れているとします。1秒間に流れる水は、断面積を底面、長さ \qty{2}{m} を高さとする円柱の形をしています（柱の体積）。内側の直径が \qty{13}{mm}（半径 \qty{0.0065}{m}）なら、断面積は $3.14\times0.0065^2=\qty{0.000133}{m^2}$。1秒に流れる水は
\bkivfline{$0.000133\times2=\qty{0.000266}{m^3}$\quad で、約 \qty{0.27}{L}}
1分では約 \qty{16}{L} です。流れる量（流量）＝断面積×速さ は、水道・空調・血管の血液の流れなど、流れを計算するときの基本の式です。
\end{minipage}\hfill
\begin{minipage}[c]{0.34\linewidth}\centering
\begin{tikzpicture}[x=1mm, y=1mm]
  \fill[bkGray!25] (0,-7) rectangle (44,7);
  \fill[bkBlue!30] (10,-6) rectangle (36,6);
  \draw[bkInk!70, line width=0.8pt] (0,7) -- (44,7) (0,-7) -- (44,-7);
  \filldraw[fill=bkBlue!45, draw=bkBlue!70!black, line width=0.6pt] (10,0) ellipse (2 and 6);
  \draw[bkBlue!70!black, line width=0.6pt] (36,6) arc (90:-90:2 and 6);
  \draw[bk arrow, bkBlue!80!black, line width=0.8pt] (14,0) -- (32,0);
  \node[font=\sffamily\footnotesize, anchor=north] at (23,-8) {1秒に進む長さ \qty{2}{m}};
  \node[font=\sffamily\footnotesize, anchor=south] at (10,7.5) {断面積};
\end{tikzpicture}
\end{minipage}

\end{butsuri}

\begin{butsuri}
\par\addvspace{3pt}\textbf{ボルトの太さと強さ}\par\noindent
\begin{minipage}[c]{0.56\linewidth}
ボルトや鉄の棒が、引っぱる力にどれだけたえられるかは、断面積に比例します。直径 \qty{10}{mm} の鉄の棒の断面積は $3.14\times5^2=\qty{78.5}{mm^2}$。直径を2倍の \qty{20}{mm} にすると、断面積は $2^2=4$ 倍の \qty{314}{mm^2} になり、4倍の力にたえられます（相似比2なら面積比4、4-3）。橋や建物の設計では、部品の断面積を計算して、必要な太さを決めます。
\end{minipage}\hfill
\begin{minipage}[c]{0.4\linewidth}\centering
\begin{tikzpicture}[x=1mm, y=1mm]
  % ボルト（横から）
  \filldraw[fill=bkGray!45, draw=bkInk!70, line width=0.6pt] (0,20) -- (4,24) -- (12,24) -- (16,20) -- (12,16) -- (4,16) -- cycle;
  \fill[bkGray!30] (5,-4) rectangle (11,16);
  \foreach \yy in {-3,-1,...,13} {\draw[bkInk!50, line width=0.4pt] (5,\yy) -- (11,\yy+1.2);}
  \draw[bkInk!70, line width=0.6pt] (5,-4) rectangle (11,16);
  \draw[bkRed, line width=0.8pt, dashed] (2,6) -- (14,6);
  \node[font=\sffamily\footnotesize, anchor=north] at (8,-5) {ボルト};
  % 断面（直径10と20）
  \filldraw[fill=bkBlue!25, draw=bkBlue!70!black, line width=0.7pt] (25,10) circle (3.2);
  \node[font=\sffamily\footnotesize, anchor=north] at (25,5.5) {\qty{10}{mm}};
  \filldraw[fill=bkBlue!25, draw=bkBlue!70!black, line width=0.7pt] (42,10) circle (6.4);
  \node[font=\sffamily\footnotesize, anchor=north] at (42,2.3) {\qty{20}{mm}};
  \node[font=\sffamily\footnotesize, anchor=south, align=center] at (33.5,17.5) {断面積は4倍};
\end{tikzpicture}
\end{minipage}

\par\addvspace{6pt}\textbf{ガスタンクが球の形をしているわけ}\par\noindent
\begin{minipage}[c]{0.66\linewidth}
同じ体積なら、球は、ほかのどんな形よりも表面積が小さくなります。体積 \qty{1}{m^3} の立方体の表面積は \qty{6}{m^2} ですが、体積 \qty{1}{m^3} の球（半径約 \qty{0.62}{m}）の表面積は約 \qty{4.84}{m^2} で、約2割少なくすみます。表面積が小さいと、タンクを作る鉄板が少なくてすみ、中の圧力も、表面に均等にかかります。都市ガスをためる大きな球形タンクには、直径 \qty{30}{m}（体積約 \qty{14000}{m^3}）ほどのものもあります。
\end{minipage}\hfill
\begin{minipage}[c]{0.3\linewidth}\centering
\begin{tikzpicture}[x=1mm, y=1mm]
  \shade[ball color=bkGray!40] (0,16) circle (15);
  \foreach \xx in {-11,-5,5,11} {\draw[bkInk!70, line width=1.2pt] (\xx,0) -- (\xx*0.85,9);}
  \draw[bkInk!60, line width=0.6pt] (-15,16) arc (180:360:15 and 4);
  \draw[bkGreen!60!black, line width=1pt] (-19,0) -- (19,0);
  \node[font=\sffamily\footnotesize, anchor=south] at (0,31.5) {球形ガスタンク};
\end{tikzpicture}
\end{minipage}

\vspace{4pt}
\begin{center}
\begin{tikzpicture}[x={(1mm,0mm)}, y={(0.45mm,0.35mm)}, z={(0mm,1mm)}]
  \bkivfbox{0}{0}{0}{16}{16}{16}{bkOrange}
  \node[font=\sffamily\footnotesize, anchor=north, align=center] at (11,0,-2) {体積 \qty{1}{m^3} の立方体\\表面積 \qty{6}{m^2}};
  \shade[ball color=bkBlue!35] (68,0,10) circle (9.9);
  \node[font=\sffamily\footnotesize, anchor=north, align=center] at (68,0,-2) {体積 \qty{1}{m^3} の球\\表面積 約 \qty{4.84}{m^2}};
\end{tikzpicture}
\end{center}
\end{butsuri}


%% ---------------- 章末問題 ----------------
\bkendproblems
\bkproblemgroup{基本}{1}

\bkproblem 次の面積を求めましょう（\bkreidairef{1}）。\quad（1）底辺 \qty{10}{cm}、高さ \qty{6}{cm} の平行四辺形\quad（2）底辺 \qty{7}{cm}、高さ \qty{8}{cm} の三角形
\bkpans{（1）\nobreak \qty{60}{cm^2}。（2）\nobreak \qty{28}{cm^2}}

\bkproblem （1）上底 \qty{5}{cm}、下底 \qty{9}{cm}、高さ \qty{4}{cm} の台形と、（2）対角線が \qty{12}{cm} と \qty{9}{cm} のひし形の面積を求めましょう（\bkreidairef{2}）。
\bkpans{（1）\nobreak \qty{28}{cm^2}。（2）\nobreak \qty{54}{cm^2}}

\bkproblem （1）半径 \qty{4}{cm} の円と、（2）半径 \qty{12}{cm}、中心角30°のおうぎ形の面積を求めましょう（\bkreidairef{4}）。
\bkpans{（1）\nobreak $16\pi$\,cm$^2$（約 \qty{50.2}{cm^2}）。（2）\nobreak $144\pi\times\dfrac{30}{360}=12\pi$\,cm$^2$（約 \qty{37.7}{cm^2}）}

\bkproblem （1）底面の半径 \qty{5}{cm}、高さ \qty{8}{cm} の円柱と、（2）底面の半径 \qty{5}{cm}、高さ \qty{9}{cm} の円すいの体積を求めましょう（\bkreidairef{5}・\bkreidairef{6}）。
\bkpans{（1）\nobreak $200\pi$\,cm$^3$（約 \qty{628}{cm^3}）。（2）\nobreak $25\pi\times9\times\dfrac{1}{3}=75\pi$\,cm$^3$（約 \qty{236}{cm^3}）}

\bkproblem 半径 \qty{3}{cm} の球の体積と表面積を求めましょう（\bkreidairef{7}）。
\bkpans{体積 $\dfrac{4}{3}\pi\times27=36\pi$\,cm$^3$、表面積 $36\pi$\,cm$^2$（どちらも約113。数は同じでも、単位がちがう）}

\bkproblem 底面の半径 \qty{6}{cm}、高さ \qty{8}{cm} の円すいの表面積を求めましょう（\bkreidairef{8}）。
\bkpans{$\ell=\sqrt{36+64}=10$。$60\pi+36\pi=96\pi$\,cm$^2$（約 \qty{301}{cm^2}）}

\bkproblem 内側が、たて \qty{40}{cm}・横 \qty{80}{cm}・深さ \qty{45}{cm} の水そうに入る水は何Lですか（\bkreidairef{10}）。
\bkpans{$40\times80\times45=\qty{144000}{cm^3}$ で、\qty{144}{L}}

\bkproblemgroup{チャレンジ}{2}

\bkproblem 四角形ABCDの対角線ACは \qty{14}{cm}、BとDからACまでの垂直な長さは \qty{5}{cm} と \qty{3}{cm} です。面積を求めましょう（\bkreidairef{3}）。
\bkpans{$14\times5\div2+14\times3\div2=35+21=\qty{56}{cm^2}$}

\Needspace{14\baselineskip}
\bkproblem 直角をはさむ辺が \qty{5}{cm} と \qty{12}{cm} の直角三角形を、\qty{12}{cm} の辺を軸にして1回転させます。できる立体の体積と表面積を求めましょう（\bkreidairef{8}・\bkreidairef{9}）。\par
\begin{center}
\begin{tikzpicture}[x=1.15mm, y=1.15mm]
  \draw[bkGray, line width=0.6pt, dash dot] (0,-1.5) -- (0,14);
  \fill[bkOrange!40] (0,0) -- (5,0) -- (0,12) -- cycle;
  \draw[bkOrange!80!black, line width=0.9pt] (0,0) -- (5,0) -- (0,12) -- cycle;
  \bkivfright{(0,0)}{(5,0)}{(0,12)}
  \node[font=\sffamily\footnotesize, anchor=north] at (2.5,-0.3) {\qty{5}{cm}};
  \node[font=\sffamily\footnotesize, anchor=east] at (-0.6,6) {\qty{12}{cm}};
  \draw[bk arrow, bkGray, line width=0.7pt] (-2.2,12.8) arc (200:340:2.3 and 0.9);
\end{tikzpicture}
\end{center}
\bkpans{半径5・高さ12の円すい。体積 $25\pi\times12\times\dfrac{1}{3}=100\pi$\,cm$^3$。母線13で、表面積 $65\pi+25\pi=90\pi$\,cm$^2$}

\bkproblem 直径 \qty{30}{m} の球形のガスタンクの体積と表面積を求めましょう（$\pi=3.14$、\bkreidairef{7}）。\par
\begin{center}
\begin{tikzpicture}[x=1mm, y=1mm]
  \shade[ball color=bkGray!40] (0,11) circle (10);
  \foreach \xx in {-7,-3,3,7} {\draw[bkInk!70, line width=1pt] (\xx,0) -- (\xx*0.85,6);}
  \draw[bkGreen!60!black, line width=1pt] (-15,0) -- (15,0);
  \draw[bkInk, line width=0.6pt, <->] (-10,11) -- (10,11);
  \node[font=\sffamily\footnotesize, anchor=south, fill=white, inner sep=1pt] at (0,11.8) {\qty{30}{m}};
  \bkperson{(19,0)}{1.13}{bkBlue}
  \draw[bkInk!60, line width=0.4pt] (19.5,1.3) -- (22,3);
  \node[font=\sffamily\footnotesize, anchor=west] at (22,3) {人（同じ縮尺）};
\end{tikzpicture}
\end{center}
\bkpans{半径15。体積 $\dfrac{4}{3}\times3.14\times3375=\qty{14130}{m^3}$、表面積 $4\times3.14\times225=\qty{2826}{m^2}$}

\begin{furikaeri}
\begin{checklist}
\item 平行四辺形・三角形・台形・ひし形の面積を、長方形に直して求められる。\bkfurref{bk:4-5:sankaku}
\item 円の面積と、柱・すい・球の体積を、わけといっしょに使える。\bkfurref{bk:4-5:en}
\item 展開図から表面積を求め、体積を L に換算できる。\bkfurref{bk:4-5:tenkai}
\end{checklist}
\end{furikaeri}

\bkfukidashi{mina}{細かく分けて、並べかえて、たす。円も球も、この考えで計算できたね。第2部では、坂道の急さを角度で表すよ。}
